\documentclass[12pt,a4paper,UTF8]{ctexart}
\usepackage[a4paper,top=22mm,bottom=22mm,left=22mm,right=22mm,headheight=15pt]{geometry}
\newcommand{\Needspace}[1]{\par\begingroup\dimen0=#1\relax\dimen2=\pagegoal\advance\dimen2 by-\pagetotal\ifdim\dimen0>\dimen2\newpage\fi\endgroup}
\newcommand{\needspace}[1]{\par\begingroup\skip0=\lastskip\vskip-\skip0\dimen0=#1\relax\advance\dimen0 by\skip0\vskip\dimen0\penalty-100\vskip-\dimen0\vskip\skip0\penalty9999\endgroup}
\usepackage{amsmath,amssymb,mathtools,bm}
\usepackage{graphicx,xcolor,booktabs,longtable,array,enumitem,fancyhdr,caption,float,microtype,setspace}
\usepackage[hidelinks,unicode]{hyperref}
\usepackage[most]{tcolorbox}
\definecolor{ContestBlue}{HTML}{17365D}
\hypersetup{pdftitle={第074套},pdfauthor={}}
\setstretch{1.24}
\setlength{\parindent}{2em}
\setlength{\parskip}{0.35em}
\setlength{\emergencystretch}{2em}
\setlist{itemsep=0.35em,topsep=0.4em}
\captionsetup{font=small,labelfont=bf,justification=centering}
\setcounter{secnumdepth}{0}
\pagestyle{fancy}
\fancyhf{}
\fancyhead[L]{\small 第074套}
\fancyfoot[C]{\thepage}
\renewcommand{\headrulewidth}{0.4pt}
\ctexset{
  section={format=\Large\bfseries\color{ContestBlue},beforeskip=1.2em,afterskip=0.7em},
  subsection={format=\large\bfseries\color{ContestBlue},beforeskip=1.0em,afterskip=0.5em},
  subsubsection={format=\normalsize\bfseries,beforeskip=0.8em,afterskip=0.35em}
}
\usepackage{calc}
\newcounter{none}
\let\originalsubsection\subsection
\renewcommand{\subsection}{\needspace{6\baselineskip}\originalsubsection}
\makeatletter
\newcommand{\tightlist}{\setlength{\itemsep}{0pt}\setlength{\parskip}{0pt}}
\newcommand{\maxwidth}{\ifdim\Gin@nat@width>\linewidth\linewidth\else\Gin@nat@width\fi}
\newcommand{\maxheight}{\ifdim\Gin@nat@height>0.72\textheight 0.72\textheight\else\Gin@nat@height\fi}
\newcommand{\pandocbounded}[1]{\begingroup\setkeys{Gin}{width=\maxwidth,height=\maxheight,keepaspectratio}#1\endgroup}
\makeatother
\begin{document}
\begin{titlepage}
  \centering
  \vspace*{58mm}
  {\Huge\bfseries\color{ContestBlue} 第074套\par}
  \vfill
  \vspace*{28mm}
\end{titlepage}
\section{第 1 题}\label{ux7b2c-1-ux9898}

一根长度为 \(L=1\ \mathrm m\)
的弹性轻杆水平悬臂放置，左端固定。在右端悬挂质量为 \(1\ \mathrm{kg}\)
的重物后，杆的右端下降
\(h=1\ \mathrm{cm}\)。现将这根杆竖直放在地面上，下端固定，在上端施加竖直向下的压力
\(F\)。估算使杆开始发生弯曲失稳所需的力 \(F\)。

\begin{figure}[H]
\centering
\includegraphics[width=0.322\linewidth,height=\textheight,keepaspectratio,alt={图1-1　竖直弹性杆的受压失稳模型}]{../figures/question-01.pdf}
\caption*{图1-1　竖直弹性杆的受压失稳模型}
\end{figure}

\subsection{一、题型定位}\label{ux4e00ux9898ux578bux5b9aux4f4d}

悬臂端点的挠度已经给出了杆的弯曲刚度。受压失稳还取决于端部约束：按下端固支、上端自由的图示，应使用悬臂压杆的最低模态；原参考的圆弧估算另作说明。

\subsection{二、知识点}\label{ux4e8cux77e5ux8bc6ux70b9}

\begin{enumerate}
\def\labelenumi{\arabic{enumi}.}
\tightlist
\item
  均匀细梁小挠度理论中，弯矩与曲率满足 \(M_b=EIy''\)。
\item
  悬臂端点横向载荷 \(W\) 引起挠度 \(h=WL^3/(3EI)\)。
\item
  弯曲弹性能为 \(U_b=(EI/2)\int_0^L(y'')^2dx\)，曲率通常不均匀。
\item
  恒定竖直压力对端点缩短量做功，势能变化为 \(-F\int_0^L(y')^2dx/2\)。
\item
  固支端满足 \(y(0)=y'(0)=0\)；自由端的自然边界条件来自能量变分。
\item
  最低允许模态决定开始失稳的临界力，较高试探模态只给较大的阈值。
\end{enumerate}

\subsection{三、物理图像}\label{ux4e09ux7269ux7406ux56feux50cf}

水平端载时，弯矩从固定端的 \(WL\)
线性减小到自由端的零值。因此曲率随位置变化，不能把整个悬臂当成等曲率圆弧来精确标定刚度。

竖直受压时，取杆长坐标为 \(x\)，横向位移为
\(y(x)\)。底端位置与切线固定，上端可横移和转动，压力方向保持竖直。微小侧弯会增加弯曲能、缩短端点高度；二次能量的两项竞争决定直杆是否稳定。

\subsection{四、解题入口}\label{ux56dbux89e3ux9898ux5165ux53e3}

先积分 \(EIy''=W(L-x)\) 求出 \(h\)，得到 \(EI=WL^3/(3h)\)。再把斜率
\(q=y'\) 作为试探函数，比较 \(EI\int(q')^2dx\) 与
\(F\int q^2dx\)。固支---自由条件给出最低波数
\(\pi/(2L)\)，无需知道杆的材料和截面。

\subsection{五、解析}\label{ux4e94ux89e3ux6790}

\begin{figure}[H]
\centering
\includegraphics[width=0.644\linewidth,height=\textheight,keepaspectratio,alt={图1-2　悬臂挠度标定与固支---自由压杆最低模态}]{../figures/solution-01.pdf}
\caption*{图1-2　悬臂挠度标定与固支---自由压杆最低模态}
\end{figure}

图1-2左侧画出端载悬臂的非均匀曲率，右侧画出满足底端固支的失稳模态。推导中的边界条件与这两个图一致。

\needspace{3\baselineskip}
\textbf{（一）由端载挠度求弯曲刚度}

取小挠度方向为正，固定端为 \(x=0\)。截面弯矩为 \(M_b=W(L-x)\)，故

\[EIy''=W(L-x).\]

由 \(y'(0)=0\) 积分一次，

\[EIy'=W(Lx-x^2/2).\]

由 \(y(0)=0\) 再积分，

\[EIy=W(Lx^2/2-x^3/6).\]

自由端挠度因此为

\[h=y(L)=\frac{WL^3}{3EI},\qquad
\boxed{EI=\frac{WL^3}{3h}}.\]

题设
\(W\approx10\ \mathrm N\)、\(L=1\ \mathrm m\)、\(h=0.01\ \mathrm m\)，得到
\(EI\approx333.3\ \mathrm{N\,m^2}\)。线弹性缓慢加载的弹性能仍为
\(Wh/2\)；将上述真实曲率代入能量积分也会得到这一结果。

\needspace{3\baselineskip}
\textbf{（二）压力势能和二次稳定性判据}

微小弯曲后，杆的端点高度缩短量为

\[\Delta=\frac12\int_0^L(y')^2dx+O(y'^4).\]

所以二次势能为

\[\Pi_2=\frac12\int_0^L\left[EI(y'')^2-F(y')^2\right]dx.\]

令 \(q=y'\)，固定端要求 \(q(0)=0\)，自由端不受弯矩要求
\(q'(L)=0\)。临界时的驻值方程为

\[EIq''+Fq=0.\]

其解 \(q=A\sin(kx)\)，其中 \(k^2=F/EI\)。自由端条件给出
\(\cos(kL)=0\)，允许波数

\[k_jL=(j+1/2)\pi,\qquad j=0,1,\ldots.\]

最低波数首先失稳，故

\[\boxed{F_c=\frac{\pi^2EI}{4L^2}
=\frac{\pi^2WL}{12h}\approx8.22\times10^2\ \mathrm N.}\]

\needspace{3\baselineskip}
\textbf{（三）直接代入最低模态核验}

取

\[y(x)=A[1-\cos(\pi x/(2L))].\]

它满足底端的位移、斜率条件及顶端零弯矩条件。两个积分为

\[\int_0^L(y')^2dx=\frac{A^2\pi^2}{8L},\qquad
\int_0^L(y'')^2dx=\frac{A^2\pi^4}{32L^3}.\]

令 \(\Pi_2=0\) 再次得到 \(F_c=\pi^2EI/(4L^2)\)。自由端横向自然条件是
\(EIy^{(3)}(L)+Fy'(L)=0\)，该模态也满足它。

\needspace{3\baselineskip}
\textbf{（四）原参考解的圆弧估算为何给出另一系数}

原解把水平悬臂近似为等曲率圆弧，得到 \(R_c\approx L^2/(2h)\)；再令
\(CL/R_c^2=Wh/2\)，得到
\(C=WL^3/(8h)\)。对于上下对称的两段圆弧，端点缩短量取
\(L^3/(24r^2)\)，于是临界估算为

\[F_{\rm arc}=24C/L^2=3WL/h\approx3.0\times10^3\ \mathrm N.\]

这一估计保留了约 \(10^3\ \mathrm N\)
的量级，但悬臂标定的等曲率近似和竖直杆两端切线约束都与图示固支---自由梁不同，不能把其系数当作同一边界条件下的严格临界值。本题若只要求原参考的估算，可注明该近似；按常规梁模型应使用约
\(0.82\ \mathrm{kN}\)。

\needspace{3\baselineskip}
\textbf{（五）端部约束改变时的检查}

两端铰支的最低模态为正弦，临界力是 \(\pi^2EI/L^2\)；两端固支则为
\(4\pi^2EI/L^2\)。这些均大于本题固支---自由值。图若把顶端位置或方向额外固定，答案系数必须随之改变。

这里忽略剪切变形、杆自重和轴向压缩，适用于细长、线弹性杆，压力为保持竖直方向的保守恒力。随杆切线方向转动的随动力是另一个稳定性模型。

\newpage

\section{第 2 题}\label{ux7b2c-2-ux9898}

在光滑水平面上有三个半径均为 \(R\) 的均匀球。球1和球2的质量均为
\(m\)，球3的质量为
\(\lambda m\)。初始时球2静止叠放在球1正上方，球3以速度 \(v\)
向右运动。球与球之间的碰撞均为完全弹性碰撞，所有接触面均无摩擦。

（1）求球3与球1第一次碰撞结束瞬间两球的速度，以向右为正方向。

（2）暂不考虑球2与球3的碰撞。求 \(\lambda\) 与 \(v\)
满足什么条件时，第一次碰撞后球2会直接进入自由落体，而不是沿球1表面滑下。

（3）在满足第（2）问的条件下，球2随后与球3相碰，且碰撞瞬间两球球心连线与水平方向成
\(45^\circ\)。求这种情形可能对应的 \(\lambda\) 范围。

（4）取
\(\lambda=3\)、\(v=\sqrt{Rg}\)，并假设球3始终不离开地面。求球2此后相对地面所能达到的最大高度
\(h\)，以 \(R\) 表示。

\begin{figure}[H]
\centering
\includegraphics[width=0.631\linewidth,height=\textheight,keepaspectratio,alt={图2-1　三球的初始位置与第一次碰撞}]{../figures/question-02.pdf}
\caption*{图2-1　三球的初始位置与第一次碰撞}
\end{figure}

\subsection{一、题型定位}\label{ux4e00ux9898ux578bux5b9aux4f4d-1}

第一次正碰决定两只底球的速度，顶部球是否脱离须另验接触条件。第二次碰撞的法线由实际接触时刻确定；第（4）问球3贴地，还要把地面冲量纳入恢复系数方程。

\subsection{二、知识点}\label{ux4e8cux77e5ux8bc6ux70b9-1}

\begin{enumerate}
\def\labelenumi{\arabic{enumi}.}
\tightlist
\item
  一维完全弹性碰撞的动量守恒与机械能守恒。
\item
  光滑球面只提供沿连心线方向的法向力或碰撞冲量，不改变切向分速度。
\item
  接触是否能够维持必须由法向力非负来判断；也可比较自由轨迹与球面几何边界的局部曲率。
\item
  在匀速运动的球1参考系中，球2若脱离便作平抛运动。
\item
  两球相碰的几何条件是球心距离等于 \(2R\)。
\item
  光滑球的斜碰可沿连心线和公切线分解；仅法向相对速度发生完全弹性反转规律。
\item
  碰后球2作竖直方向的抛体运动，最高点由 \(v_y=0\) 或
  \(\Delta h=v_{y0}^2/(2g)\) 求得。
\end{enumerate}

\subsection{三、物理图像}\label{ux4e09ux7269ux7406ux56feux50cf-1}

第一次碰撞是严格的一维正碰，球2因接触点在球1正上方且没有摩擦，在极短碰撞过程中不会获得水平冲量。碰后若球1向右运动足够快，球2来不及沿球1表面``贴着下落''，二者接触立即断开；在球1随动系中，球2具有向左的水平初速度和向下的重力加速度。

球2下落到球1左侧时可能遇到落后的球3。第二次碰撞的法线沿球3指向球2，与水平成
\(45^\circ\)。这时必须把碰前速度投影到该法线上，再用质量不同的完全弹性碰撞公式；不能把两个速度的水平、竖直分量分别做一维碰撞。

\subsection{四、解题入口}\label{ux56dbux89e3ux9898ux5165ux53e3-1}

第一问先完整解出 \(v_{10}\)、\(v_{30}\)。第二问有两种等价入口：

\begin{itemize}
\tightlist
\item
  在球1参考系写球2的抛物线，并要求它在初始点附近位于半径 \(2R\)
  的禁入圆之外；
\item
  在刚要脱离的顶部写法向动力学，要求维持接触所需的法向力不为正。
\end{itemize}

前一种方法与后续几何衔接更自然。第三问先用给定的 \(45^\circ\)
几何确定球2下降的水平、竖直位移，再比较球2与球3的运动时间。第四问沿连心线做一次标准斜碰，最后只保留球2碰后竖直速度。

\subsection{五、解析}\label{ux4e94ux89e3ux6790-1}

\begin{figure}[H]
\centering
\includegraphics[width=0.692\linewidth,height=\textheight,keepaspectratio,alt={图2-2　球2的脱离判据与球2、球3的斜碰几何}]{../figures/solution-02.pdf}
\caption*{图2-2　球2的脱离判据与球2、球3的斜碰几何}
\end{figure}

图 2-2 左侧给出球 2
脱离时支持力恰为零的径向关系，右侧把球2、球3的连心法线与地面支持冲量分开。以下先用左图确定碰前速度，再按右图写球间冲量和球3贴地约束，再使用法向恢复系数。

\needspace{3\baselineskip}
\textbf{（一）第一次完全弹性正碰}

设第一次碰撞后球1、球3的速度分别为
\(v_{10}\)、\(v_{30}\)。碰前球1静止，球3速度为
\(v\)。沿水平方向动量守恒：

\[
\lambda m v
=m v_{10}+\lambda m v_{30}. \tag{1}
\]

约去 \(m\)：

\[
\lambda v=v_{10}+\lambda v_{30}. \tag{2}
\]

完全弹性碰撞还满足动能守恒：

\[
\frac12\lambda m v^2
=\frac12m v_{10}^2
+\frac12\lambda m v_{30}^2. \tag{3}
\]

约去 \(m/2\)：

\[
\lambda v^2=v_{10}^2+\lambda v_{30}^2. \tag{4}
\]

也可使用完全弹性正碰的相对速度关系：

\[
v_{10}-v_{30}=v-0=v. \tag{5}
\]

由式（5）得 \(v_{10}=v+v_{30}\)，代入式（2）：

\[
\lambda v=v+(1+\lambda)v_{30}.
\]

所以

\[
v_{30}=\frac{\lambda-1}{\lambda+1}v. \tag{6}
\]

再代回式（5）：

\[
\begin{aligned}
v_{10}
&=v+\frac{\lambda-1}{\lambda+1}v\\
&=\frac{\lambda+1+\lambda-1}{\lambda+1}v\\
&=\frac{2\lambda}{\lambda+1}v. 
\end{aligned}
\tag{7}
\]

因此

\[
\boxed{v_{30}=\frac{\lambda-1}{\lambda+1}v,\qquad
v_{10}=\frac{2\lambda}{\lambda+1}v.}
\]

若 \(\lambda>1\)，球3碰后仍向右；若 \(\lambda=1\)，球3停止；若
\(0<\lambda<1\)，球3反弹向左。第三问要让球2与球3在球1左侧相碰，最终会要求
\(\lambda>1\)。

\needspace{3\baselineskip}
\textbf{（二）球2直接脱离球1的条件}

取球1碰后的匀速参考系。以球2初始球心为坐标原点，取水平向右为 \(x\)
正方向、竖直向上为 \(y\) 正方向。碰撞结束瞬间：

\begin{itemize}
\tightlist
\item
  球1在该参考系静止；
\item
  球2在实验室中水平速度仍为零，所以相对球1的水平速度为 \(-v_{10}\)；
\item
  球2的初始竖直速度为零。
\end{itemize}

若球2立即脱离，它的自由运动方程为

\[
x=-v_{10}t,
\qquad
y=-\frac12gt^2. \tag{8}
\]

由第一式得到

\[
t=-\frac{x}{v_{10}}.
\]

代入第二式，消去时间：

\[
y=-\frac{g}{2v_{10}^2}x^2. \tag{9}
\]

记

\[
\beta=\frac{g}{2v_{10}^2},
\]

则自由轨迹为 \(y=-\beta x^2\)。

在所取坐标系中，球1球心位于 \((0,-2R)\)。为了使两个半径为 \(R\)
的球不重叠，球2球心到球1球心的距离必须满足

\[
x^2+(y+2R)^2\ge(2R)^2. \tag{10}
\]

将 \(y=-\beta x^2\) 代入：

\[
x^2+(2R-\beta x^2)^2\ge4R^2.
\]

展开平方：

\[
x^2+4R^2-4R\beta x^2+\beta^2x^4\ge4R^2.
\]

消去两边的 \(4R^2\)：

\[
\beta^2x^4+(1-4R\beta)x^2\ge0. \tag{11}
\]

在刚离开初始点时 \(x\) 是任意小量，\(x^4\) 比 \(x^2\)
高阶。要使自由轨迹一开始就在球面之外，\(x^2\) 项系数必须非负：

\[
1-4R\beta\ge0. \tag{12}
\]

代入 \(\beta=g/(2v_{10}^2)\)：

\[
1-\frac{2Rg}{v_{10}^2}\ge0,
\]

即

\[
v_{10}^2\ge2Rg. \tag{13}
\]

再代入式（7）：

\[
\left(\frac{2\lambda}{\lambda+1}v\right)^2\ge2Rg.
\]

所以球2直接进入自由落体的条件为

\[
\boxed{
\frac{4\lambda^2}{(\lambda+1)^2}v^2\ge2Rg
}
\]

或写成

\[
\boxed{
v\ge\frac{\lambda+1}{\sqrt2\lambda}\sqrt{Rg}
}. \tag{14}
\]

等号对应临界情形：自由抛物线与球面轨迹在顶部具有相同的二阶曲率。

\needspace{3\baselineskip}
\textbf{（三）由 \(45^\circ\) 碰撞几何确定 \(\lambda\) 范围}

第一次碰撞结束时，把球1球心的水平位置取为
\(x=0\)。球2球心初始位于球1正上方。球2与球3相碰时，球3仍在地面上，其球心高度为
\(R\)；球2与球3的球心连线长 \(2R\)，并与水平成 \(45^\circ\)。

因此，球2球心相对于球3球心的水平、竖直分量均为

\[
2R\cos45^\circ=\sqrt2R.
\]

球2初始球心高度为 \(3R\)，碰撞时高度为

\[
R+\sqrt2R=(1+\sqrt2)R.
\]

故球2的下降量为

\[
\Delta y
=3R-(1+\sqrt2)R
=(2-\sqrt2)R. \tag{15}
\]

球2从静止竖直下落，所以

\[
\frac12gt_c^2=(2-\sqrt2)R. \tag{16}
\]

球3第一次碰撞后从球1左侧相距 \(2R\)
处开始向右运动。到第二次碰撞时，球2位于球3右上方，二者水平间距为
\(\sqrt2R\)。球2在实验室中没有水平速度，因此球3需要向右走过

\[
2R-\sqrt2R=(2-\sqrt2)R. \tag{17}
\]

于是

\[
v_{30}t_c=(2-\sqrt2)R. \tag{18}
\]

由式（18）得

\[
t_c=\frac{(2-\sqrt2)R}{v_{30}}.
\]

代入式（16）：

\[
\frac12g\left[\frac{(2-\sqrt2)R}{v_{30}}\right]^2
=(2-\sqrt2)R.
\]

约去正数 \((2-\sqrt2)R\)：

\[
v_{30}^2=\frac{2-\sqrt2}{2}Rg. \tag{19}
\]

再用 \(v_{30}=(\lambda-1)v/(\lambda+1)\)：

\[
\left(\frac{\lambda-1}{\lambda+1}\right)^2v^2
=\frac{2-\sqrt2}{2}Rg.
\]

因此

\[
v^2
=\frac{2-\sqrt2}{2}
\left(\frac{\lambda+1}{\lambda-1}\right)^2Rg. \tag{20}
\]

要使球3碰后仍向右运动并在球2左下方追上碰撞，必须有

\[
\lambda>1. \tag{21}
\]

同时还必须满足第（2）问的脱离条件。把式（20）代入式（14）的平方形式：

\[
\frac{4\lambda^2}{(\lambda+1)^2}
\cdot
\frac{2-\sqrt2}{2}
\left(\frac{\lambda+1}{\lambda-1}\right)^2Rg
\ge2Rg.
\]

约去 \(2Rg\) 与 \((\lambda+1)^2\)，得到

\[
(2-\sqrt2)\frac{\lambda^2}{(\lambda-1)^2}\ge1. \tag{22}
\]

在 \(\lambda>1\) 下两边开正平方根：

\[
\sqrt{2-\sqrt2}\,\lambda\ge\lambda-1.
\]

移项：

\[
1\ge\lambda\left(1-\sqrt{2-\sqrt2}\right).
\]

所以

\[
\lambda\le
\frac{1}{1-\sqrt{2-\sqrt2}}
\simeq4.262. \tag{23}
\]

结合式（21）：

\[
\boxed{
1<\lambda\le
\frac{1}{1-\sqrt{2-\sqrt2}}
\simeq4.26
}. \tag{24}
\]

\needspace{3\baselineskip}
\textbf{（四）第（4）问必须重新求碰撞几何}

第（3）问的 \(45^\circ\) 规定会锁定 \(v\) 与 \(\lambda\) 的关系。代入
\(\lambda=3\)，需有 \(v^2=2(2-\sqrt2)Rg\)，与第（4）问 \(v^2=Rg\)
不同。因此本问不能继续使用 \(45^\circ\)。

以 \(R\) 为长度单位、\(\sqrt{R/g}\) 为时间单位，记无量纲时间为
\(\tau\)。第一次碰后球1速度为 \(1.5\sqrt{Rg}\)，球3速度为
\(0.5\sqrt{Rg}\)。球2直接脱离，实验室坐标为

\[\frac{x_2}R=0,\quad \frac{y_2}R=3-\frac{\tau^2}{2},\qquad
\frac{x_3}R=-2+\frac\tau2,\quad \frac{y_3}R=1.\]

球心距离首次等于 \(2R\)，故

\[\left(2-\frac\tau2\right)^2+
\left(2-\frac{\tau^2}{2}\right)^2=4.\]

即

\[\tau^4-7\tau^2-8\tau+16=0.\]

取最早的正接近根 \(\tau_c=1.11073996\)。法线从球3指向球2，

\[\boldsymbol n=(n_x,n_y)=(0.72231501,0.69156419),\]

夹角约 \(43.754^\circ\)，接触高度 \(y_c/R=2.38312837\)。

\needspace{3\baselineskip}
\textbf{（五）球3保持贴地的碰撞冲量}

以 \(\sqrt{Rg}\) 为速度单位，碰前球2为 \((0,-\tau_c)\)，球3为
\((1/2,0)\)。对球2的冲量是
\(J\boldsymbol n\)；球3受到反向冲量及地面向上的冲量
\(J n_y\)，因而竖直速度保持零。设 \(j=J/(m\sqrt{Rg})\)，碰后

\[\boldsymbol u_2^+=(jn_x,-\tau_c+jn_y),\qquad
\boldsymbol u_3^+=(1/2-jn_x/3,0).\]

恢复系数为1，法向相对速度反号，所以

\[j\left(1+\frac{n_x^2}{3}\right)
=2\left(\frac{n_x}{2}+\tau_cn_y\right).\]

注意分母是 \(1+n_x^2/3\)；使用 \(1+1/3\)
会让球3获得向下速度，与贴地条件矛盾。得到

\[j=1.92400202,\qquad u_{2y}^+=0.219830933>0.\]

\needspace{3\baselineskip}
\textbf{（六）最高点、守恒与后续接触检查}

球2再上升 \(R(u_{2y}^+)^2/2\)，故

\[\boxed{h_{\max}/R=2.38312837+0.024162820\approx2.40729119\approx2.41.}\]

这是球心高度；若指球的最高表面点，还要加一个半径。碰撞的水平总动量守恒，竖直总动量因地面冲量改变，动能仍守恒，因为理想固定地面在碰撞中不做功。数值代回两球动能可以核对上述
\(j\)。

首次碰撞后球1与球2的距离始终增大，球1与球3也已分离；在第二次碰撞到球2达到此峰点的时间内，球2未再次接触球1或球3。原稿把45°几何和自由两球斜碰公式同时沿用，所报
\(2.41966R\) 不是本问给定参数下的结果。

以上最高点指本问第二次碰撞后的第一次抛体峰值；若再考虑落地、反弹和长期多次碰撞，还需题目给出球---地面的恢复规则。

\needspace{3\baselineskip}
\textbf{第（3）问的范围与首次接触条件}

在45°规定下，\(v_{30}=\sqrt{(2-\sqrt2)Rg/2}>0\)；这个速度由接触几何确定，而不是可随意指定。与脱离判据联立后得到
\(1<\lambda\leq4.26197\)。沿该自由轨迹，球2---球3距离在到达45°之前单调减小，所取确为首次接触。

顶部脱离条件等号时，距离平方减去 \(4R^2\)
的二阶项消失，但四阶项仍为正，所以自由轨迹仍位于禁入圆之外。接触支持力只能推不能拉，是判断等号可以纳入的依据。

\newpage

\section{第 3 题}\label{ux7b2c-3-ux9898}

两根相距 \(l\) 的无限长带电直线均平行于 \(z\)
轴固定放置，线电荷密度分别为 \(+\lambda\) 和 \(-\lambda\)。在
\(x\)---\(y\)
平面内，负线电荷位于左侧，正线电荷位于右侧，连接两线电荷的主轴正方向向右。取两线电荷中点为极坐标原点。

（1）当 \(r\gg l\) 时，求点 \(A(r,\theta)\) 的电势。

（2）求该点的径向电场分量 \(E_r\) 和角向电场分量
\(E_\theta\)，并说明等势线与电场线的形状。

（3）移去线电荷 \(-\lambda\)，并将线电荷 \(+\lambda\) 平移到原点。在
\((0,r_0)\) 处放置质量为 \(m\)、电荷量为 \(-q\) 的粒子，给它沿 \(+x\)
方向的初速度

\[
v_0=\sqrt{\frac{2\lambda q}{\pi\varepsilon_0m}}.
\]

求：

I. 粒子的径向速度下一次等于零时，它到原点的距离 \(r_1\)；

\begin{enumerate}
\def\labelenumi{\Roman{enumi}.}
\setcounter{enumi}{1}
\tightlist
\item
  从初始时刻到该时刻，粒子绕原点转过的角度 \(\theta_0\)。
\end{enumerate}

\begin{figure}[H]
\centering
\includegraphics[width=0.628\linewidth,height=\textheight,keepaspectratio,alt={图3-1　异号线电荷在横截面内的位置}]{../figures/question-03.pdf}
\caption*{图3-1　异号线电荷在横截面内的位置}
\end{figure}

\subsection{一、题型定位}\label{ux4e00ux9898ux578bux5b9aux4f4d-2}

异号线电荷的对数势相减，远区先得到二维偶极势。移去负线后，粒子进入吸引的对数中心势；先求两个径向转向点，再用角动量消去时间计算扫角。

\subsection{二、知识点}\label{ux4e8cux77e5ux8bc6ux70b9-2}

\begin{enumerate}
\def\labelenumi{\arabic{enumi}.}
\tightlist
\item
  无限长线电荷的径向电场大小为
  \(E=\lambda/(2\pi\varepsilon_0\rho)=2k\lambda/\rho\)，其中
  \(k=1/(4\pi\varepsilon_0)\)。
\item
  单根正线电荷的电势可写成
  \(\Phi=-2k\lambda\ln\rho+\text{常量}\)；只有电势差有确定意义。
\item
  当 \(r\gg l\) 时，对两线电荷距离之比作一阶展开，单极项相消，剩余偶极项
  \(\propto l\cos\theta/r\)。
\item
  极坐标中的电场分量：\(E_r=-\partial\Phi/\partial r\)，\(E_\theta=-(1/r)\partial\Phi/\partial\theta\)。
\item
  等势线与电场线处处正交；平面极坐标中的场线满足
  \(dr/(r\,d\theta)=E_r/E_\theta\)。
\item
  中心力运动的角动量守恒 \(J=mr^2\dot\theta\)。
\item
  对数势 \(U(r)=2k\lambda q\ln r\) 与有效势方法。
\item
  变量代换 \(u=1/r\)、\(A=r_0u\) 将径向方程转成无量纲轨道方程。
\end{enumerate}

\subsection{三、物理图像}\label{ux4e09ux7269ux7406ux56feux50cf-2}

一正一负两根线电荷的总线电荷量为零，所以远处看不到 \(1/r\)
型的二维单极场，第一非零项由``电荷强度乘分离距离'' \(\lambda l\)
决定。沿主轴正方向电势为正，沿主轴负方向电势为负，垂直平分线上电势为零。

第三问中负粒子受到指向正线电荷的吸引力。给定初速度大于同半径圆周运动速度，粒子先向外运动；角动量使它不会沿径向逃走。到达
\(r_1\) 时径向速度为零，但切向速度不为零，此点是轨道的远拱点。

\subsection{四、解题入口}\label{ux56dbux89e3ux9898ux5165ux53e3-2}

前两问从``正线电荷到场点的距离 \(\rho_+\)''与``负线电荷到场点的距离
\(\rho_-\)''写起，先把总势写成对数比，再展开；这样电势常量自动消去。

第三问先用能量和角动量求远拱点比值
\(\beta=r_0/r_1\)。求转角时，不直接联立
\(r(t)\)、\(\theta(t)\)，而是用角动量把时间消去，建立
\(A(\theta)=r_0/r\) 的二阶方程；再把二阶方程乘以 \(A'\)
积分一次，得到可数值计算的角积分。

\subsection{五、解析}\label{ux4e94ux89e3ux6790-2}

\begin{figure}[H]
\centering
\includegraphics[width=0.9\linewidth,height=\textheight,keepaspectratio,alt={图3-2　线偶极子的场线结构与对数中心势轨道}]{../figures/solution-03.pdf}
\caption*{图3-2　线偶极子的场线结构与对数中心势轨道}
\end{figure}

图 3-2 中场点到正、负线电荷的距离分别记为
\(\rho_+\)、\(\rho_-\)；远场展开与中心势轨道都由这两个距离出发。后文的极坐标
\(r,\theta\) 和拱点半径 \(r_0,r_1\) 均沿图中定义使用。

\needspace{3\baselineskip}
\textbf{（一）远区电势}

令

\[
k=\frac{1}{4\pi\varepsilon_0}.
\]

单根线电荷产生的电场为

\[
E(\rho)=\frac{\lambda}{2\pi\varepsilon_0\rho}
=\frac{2k\lambda}{\rho}. \tag{1}
\]

对正线电荷，取任意参考半径后，其电势可写为

\[
\Phi_+(\rho_+)=-2k\lambda\ln\rho_++C_+.
\]

对负线电荷，电势为

\[
\Phi_-(\rho_-)=+2k\lambda\ln\rho_-+C_-.
\]

总电势中常量可统一选取，使

\[
\Phi
=2k\lambda\ln\frac{\rho_-}{\rho_+}. \tag{2}
\]

正、负线电荷分别位于 \(x=+l/2\) 与 \(x=-l/2\)。点 \(A\) 的直角坐标为

\[
x=r\cos\theta,
\qquad
y=r\sin\theta.
\]

故

\[
\begin{aligned}
\rho_+^2
&=\left(r\cos\theta-\frac l2\right)^2+r^2\sin^2\theta\\
&=r^2-rl\cos\theta+\frac{l^2}{4},\\[4pt]
\rho_-^2
&=\left(r\cos\theta+\frac l2\right)^2+r^2\sin^2\theta\\
&=r^2+rl\cos\theta+\frac{l^2}{4}. 
\end{aligned}
\tag{3}
\]

当 \(r\gg l\) 时，保留 \(l/r\) 的一阶项：

\[
\begin{aligned}
\rho_+
&=r\sqrt{1-\frac lr\cos\theta+\frac{l^2}{4r^2}}
\simeq r\left(1-\frac{l\cos\theta}{2r}\right),\\
\rho_-
&=r\sqrt{1+\frac lr\cos\theta+\frac{l^2}{4r^2}}
\simeq r\left(1+\frac{l\cos\theta}{2r}\right). 
\end{aligned}
\tag{4}
\]

因此

\[
\frac{\rho_-}{\rho_+}
\simeq
\frac{1+\dfrac{l\cos\theta}{2r}}
{1-\dfrac{l\cos\theta}{2r}}.
\]

令 \(\delta=l\cos\theta/(2r)\)。由于 \(|\delta|\ll1\)，

\[
\ln\frac{1+\delta}{1-\delta}
=\left(\delta-\frac{\delta^2}{2}+\cdots\right)
-\left(-\delta-\frac{\delta^2}{2}+\cdots\right)
\simeq2\delta.
\]

所以

\[
\ln\frac{\rho_-}{\rho_+}
\simeq\frac{l\cos\theta}{r}. \tag{5}
\]

代入式（2）：

\[
\boxed{
\Phi(r,\theta)
=\frac{2k\lambda l\cos\theta}{r}
=\frac{\lambda l\cos\theta}{2\pi\varepsilon_0r}
}. \tag{6}
\]

沿 \(+x\) 轴有 \(\theta=0\)，电势为正；沿 \(-x\) 轴有
\(\theta=\pi\)，电势为负，符号与电荷位置相符。

\needspace{3\baselineskip}
\textbf{（二）电场分量、等势线与电场线}

为简化书写，记

\[
p_\ell=2k\lambda l.
\]

则

\[
\Phi=\frac{p_\ell\cos\theta}{r}. \tag{7}
\]

径向分量为

\[
\begin{aligned}
E_r
&=-\frac{\partial\Phi}{\partial r}\\
&=-\frac{\partial}{\partial r}
\left(p_\ell\cos\theta\,r^{-1}\right)\\
&=\frac{p_\ell\cos\theta}{r^2}. 
\end{aligned}
\tag{8}
\]

角向分量为

\[
\begin{aligned}
E_\theta
&=-\frac1r\frac{\partial\Phi}{\partial\theta}\\
&=-\frac1r\left(-\frac{p_\ell\sin\theta}{r}\right)\\
&=\frac{p_\ell\sin\theta}{r^2}. 
\end{aligned}
\tag{9}
\]

即

\[
\boxed{
E_r=\frac{2k\lambda l\cos\theta}{r^2},
\qquad
E_\theta=\frac{2k\lambda l\sin\theta}{r^2}
}. \tag{10}
\]

接着求等势线。令 \(\Phi=\Phi_0\)，由式（7）有

\[
r=\frac{p_\ell}{\Phi_0}\cos\theta. \tag{11}
\]

记 \(a=p_\ell/\Phi_0\)，再用 \(r^2=x^2+y^2\)、\(r\cos\theta=x\)：

\[
r^2=ax,
\]

即

\[
x^2+y^2=ax.
\]

配方：

\[
\left(x-\frac a2\right)^2+y^2
=\left(\frac a2\right)^2. \tag{12}
\]

所以每一条非零等势线都是圆心在 \(x\) 轴上、经过原点的圆；\(y\)
轴本身对应零电势线。

电场线的切线方向与电场方向一致，极坐标轨迹满足

\[
\frac{dr}{r\,d\theta}=\frac{E_r}{E_\theta}.
\]

代入式（10）：

\[
\frac{dr}{r\,d\theta}=\cot\theta.
\]

分离变量：

\[
\frac{dr}{r}=\cot\theta\,d\theta.
\]

积分得

\[
\ln r=\ln|\sin\theta|+\ln b,
\]

故

\[
r=b\sin\theta. \tag{13}
\]

使用 \(r\sin\theta=y\)：

\[
r^2=by,
\]

即

\[
x^2+y^2=by.
\]

配方：

\[
x^2+\left(y-\frac b2\right)^2
=\left(\frac b2\right)^2. \tag{14}
\]

所以电场线也是过原点的一族圆，但圆心位于 \(y\) 轴上。两族圆处处正交。

\needspace{3\baselineskip}
\textbf{（三）远拱点距离 \(r_1\)}

移去负线电荷后，正线电荷位于原点。它的电势可写为

\[
\Phi_\lambda(r)=-2k\lambda\ln r+\text{常量}.
\]

粒子电荷为 \(-q\)，所以其电势能为

\[
U(r)=(-q)\Phi_\lambda(r)
=2k\lambda q\ln r+\text{常量}. \tag{15}
\]

由 \(-dU/dr=-2k\lambda q/r\) 可检查，力确实沿径向向内。

初始位置为 \((0,r_0)\)，初速度沿
\(+x\)，与半径垂直。因此初始径向速度为零，角动量大小为

\[
J=mv_0r_0. \tag{16}
\]

到达下一径向速度为零的位置 \(r_1\) 时，速度全部为切向。由角动量守恒：

\[
mv_0r_0=mv_1r_1.
\]

故

\[
v_1=v_0\frac{r_0}{r_1}. \tag{17}
\]

能量守恒为

\[
\frac12mv_0^2+2k\lambda q\ln r_0
=\frac12mv_1^2+2k\lambda q\ln r_1. \tag{18}
\]

定义无量纲比值

\[
\beta=\frac{r_0}{r_1}. \tag{19}
\]

因为粒子先向外运动，\(r_1>r_0\)，故 \(0<\beta<1\)。由式（17）有

\[
v_1=\beta v_0. \tag{20}
\]

把式（20）代入式（18），并把对数项移到右侧：

\[
\frac12mv_0^2(1-\beta^2)
=2k\lambda q\ln\frac{r_1}{r_0}.
\]

由于 \(r_1/r_0=1/\beta\)：

\[
1-\beta^2
=\frac{4k\lambda q}{mv_0^2}\ln\frac1\beta. \tag{21}
\]

题给

\[
v_0^2=\frac{2\lambda q}{\pi\varepsilon_0m}.
\]

而

\[
4k=\frac1{\pi\varepsilon_0},
\]

所以

\[
\frac{4k\lambda q}{mv_0^2}
=\frac{\lambda q/(\pi\varepsilon_0m)}
{2\lambda q/(\pi\varepsilon_0m)}
=\frac12. \tag{22}
\]

式（21）化为

\[
1-\beta^2
=\frac12\ln\frac1\beta
=-\frac12\ln\beta. \tag{23}
\]

即

\[
f(\beta)=1-\beta^2+\frac12\ln\beta=0. \tag{24}
\]

\(\beta=1\) 对应初始近拱点，不是``下一次''径向速度为零的位置。另一个位于
\((0,1)\) 内的根数值为

\[
\beta=0.1408098053\ldots \tag{25}
\]

因此

\[
\boxed{
r_1=\frac{r_0}{\beta}
=7.10178\,r_0
\simeq7.10r_0
}. \tag{26}
\]

\needspace{3\baselineskip}
\textbf{（四）建立轨道方程}

令

\[
u(\theta)=\frac1r.
\]

角动量守恒给出

\[
J=mr^2\dot\theta,
\]

因此

\[
\dot\theta=\frac{J}{mr^2}=\frac{J}{m}u^2. \tag{27}
\]

对 \(r=1/u\) 求时间导数：

\[
\begin{aligned}
\dot r
&=\frac{dr}{d\theta}\dot\theta\\
&=\left(-\frac1{u^2}\frac{du}{d\theta}\right)
\left(\frac Jm u^2\right)\\
&=-\frac Jm\frac{du}{d\theta}. 
\end{aligned}
\tag{28}
\]

再求一次时间导数：

\[
\begin{aligned}
\ddot r
&=-\frac Jm\frac{d}{dt}\left(\frac{du}{d\theta}\right)\\
&=-\frac Jm\frac{d^2u}{d\theta^2}\dot\theta\\
&=-\left(\frac Jm\right)^2u^2\frac{d^2u}{d\theta^2}. 
\end{aligned}
\tag{29}
\]

极坐标径向加速度为

\[
a_r=\ddot r-r\dot\theta^2.
\]

其中

\[
r\dot\theta^2
=\frac1u\left(\frac Jm u^2\right)^2
=\left(\frac Jm\right)^2u^3. \tag{30}
\]

所以

\[
a_r
=-\left(\frac Jm\right)^2u^2
\left(\frac{d^2u}{d\theta^2}+u\right). \tag{31}
\]

径向力为

\[
F_r=-\frac{2k\lambda q}{r}=-2k\lambda q\,u. \tag{32}
\]

由 \(ma_r=F_r\)：

\[
-\frac{J^2}{m}u^2(u''+u)
=-2k\lambda q\,u.
\]

约去 \(-u\)：

\[
\frac{J^2}{m}u(u''+u)=2k\lambda q. \tag{33}
\]

由 \(J=mv_0r_0\)，有

\[
\frac{2mk\lambda q}{J^2}
=\frac{2k\lambda q}{mv_0^2r_0^2}.
\]

又由题给 \(v_0^2=8k\lambda q/m\)，所以

\[
\frac{2k\lambda q}{mv_0^2r_0^2}
=\frac1{4r_0^2}. \tag{34}
\]

式（33）变为

\[
u(u''+u)=\frac1{4r_0^2}. \tag{35}
\]

定义

\[
A=r_0u=\frac{r_0}{r}. \tag{36}
\]

由于 \(u=A/r_0\)、\(u''=A''/r_0\)，式（35）化为

\[
A(A''+A)=\frac14. \tag{37}
\]

即

\[
A''=\frac1{4A}-A. \tag{38}
\]

\needspace{3\baselineskip}
\textbf{（五）把轨道方程积分一次并求转角}

将式（38）两边乘以 \(A'\)：

\[
A''A'=\left(\frac1{4A}-A\right)A'.
\]

左边是 \((A'^2/2)'\)，右边分别可积：

\[
\frac12A'^2
=\frac14\ln A-\frac12A^2+C. \tag{39}
\]

初始时 \(r=r_0\)，故 \(A=1\)；初始径向速度为零，由式（28）知
\(u'=0\)，因而 \(A'=0\)。代入式（39）：

\[
0=0-\frac12+C,
\]

所以 \(C=1/2\)。于是

\[
A'^2=1-A^2+\frac12\ln A. \tag{40}
\]

粒子从 \(r_0\) 向外运动时，\(A=r_0/r\) 从 \(1\) 减小到 \(\beta\)，因此
\(A'<0\)：

\[
\frac{dA}{d\theta}
=-\sqrt{1-A^2+\frac12\ln A}. \tag{41}
\]

移项：

\[
d\theta
=-\frac{dA}{\sqrt{1-A^2+\dfrac12\ln A}}.
\]

从 \(A=1\) 积到 \(A=\beta\)：

\[
\theta_0
=\int_{\beta}^{1}
\frac{dA}{\sqrt{1-A^2+\dfrac12\ln A}}. \tag{42}
\]

上下限处根号均趋于零，但都属于可积的平方根型端点。代入
\(\beta=0.1408098053\ldots\) 数值积分，得到

\[
\boxed{
\theta_0=2.07814\ \mathrm{rad}
\simeq119.1^\circ
}. \tag{43}
\]

\needspace{3\baselineskip}
\textbf{角方向、远区图和独立轨道检验}

初始粒子位于 \((0,r_0)\)，速度沿
\(+x\)，所以有符号角动量为负，极角随时间减小。正文 \(J\) 及 \(\theta_0\)
用其大小，计算轨道时可用顺时针扫角 \(\varphi\geq0\)，实际极角为
\(\pi/2-\varphi\)。两次求导所得 Binet 方程不受这一改记影响，但不能同时把
\(J>0\)、\(\dot\theta>0\) 与初始速度箭头画在同一图中。

\Needspace{18\baselineskip}

用长度 \(r_0\)、时间 \(r_0/v_0\) 归一化，直接积分笛卡尔方程

\[\ddot{\boldsymbol r}=-\frac{\boldsymbol r}{4r^2},\qquad
\boldsymbol r(0)=(0,1),\quad \dot{\boldsymbol r}(0)=(1,0).\]

在 \(\boldsymbol r\cdot\dot{\boldsymbol r}\) 从正变负的首次事件读出
\(r_1/r_0\approx7.10178\)、扫角
\(2.07814\)，与正文根和端点积分一致。图中的远拱点必须离中心约
\(7.1r_0\)，不能画在半径 \(r_0\) 的圆上。

等势圆、场线圆来自 \(r\gg l\)
的近似势；它们靠近原点的部分超出远区适用范围。图中以灰色小区遮出这一范围，并标明两族的含义。

\newpage

\section{第 4 题}\label{ux7b2c-4-ux9898}

边长为 \(L\) 的正方形平行板电容器，两极板间距为
\(d\)。先把电容器接到电压为 \(V\)
的电池上，充电完成后断开电池；再将电容器竖直插入相对介电常数为
\(\varepsilon_r\)、密度为 \(\rho\)
的液体中。稳定时，液体在两极板间所占的竖直长度为
\(L/2\)。忽略其他外力影响，求：

（1）此时电容器的电容；

（2）两极板间的电场强度；

（3）金属极板自由电荷的面密度；

（4）极板内外液面的高度差 \(h\)。

\begin{figure}[H]
\centering
\includegraphics[width=0.617\linewidth,height=\textheight,keepaspectratio,alt={图4-1　方形平行板电容器的竖直浸液状态}]{../figures/question-04.pdf}
\caption*{图4-1　方形平行板电容器的竖直浸液状态}
\end{figure}

\subsection{一、题型定位}\label{ux4e00ux9898ux578bux5b9aux4f4d-3}

断开电池后自由电荷保持不变。空气区与液体区跨接同一对极板，先并联求电容和场强，再从一般浸入长度的电场能求液体上升力。

\subsection{二、知识点}\label{ux4e8cux77e5ux8bc6ux70b9-3}

\begin{enumerate}
\def\labelenumi{\arabic{enumi}.}
\tightlist
\item
  沿极板面积方向分区的介质等效为并联电容器；沿板间距方向分层才等效为串联。
\item
  断开电池后 \(Q\) 不变，电容器电压为 \(V'=Q/C\)。
\item
  均匀介质中 \(\bm D=\varepsilon\bm E\)，金属表面自由电荷密度满足
  \(\sigma_f=\bm D\cdot\bm n\)。
\item
  固定电荷条件下电场能为 \(U_E=Q^2/(2C)\)。
\item
  广义力由 \(F_x=-dU_E/dx\)
  求得；介质进入使电容增大、电场能降低，因此电力把液体向上吸入。
\item
  液面高度差 \(h\) 对应静水压差
  \(\rho gh\)，作用面积为极板宽度与板间距的乘积 \(Ld\)。
\end{enumerate}

\subsection{三、物理图像}\label{ux4e09ux7269ux7406ux56feux50cf-3}

液体进入下半部分后，电场线仍大体垂直极板。上半部分面积为
\(L^2/2\)，介质为空气；下半部分面积也为 \(L^2/2\)，介质的介电常数为
\(\varepsilon_r\varepsilon_0\)。二者跨接同一对极板，电压相同，所以是并联关系。

断电后，总自由电荷不变，但它会在金属表面重新分布。浸液部分的
\(D=\varepsilon E\) 较大，因此其自由电荷面密度比空气部分大
\(\varepsilon_r\)
倍。液体受到向上的电场拉力，直至它与由内外液面差产生的向下静水压力合力平衡。

\subsection{四、解题入口}\label{ux56dbux89e3ux9898ux5165ux53e3-3}

先回到充电完成、尚未插液的状态，求出被``锁定''的总电荷
\(Q=C_0V\)。随后对任意浸入长度 \(x\) 写出 \(C(x)\)，这样前 3 问只需令
\(x=L/2\)，第 4 问则可对一般式求导得到电场拉力。直接只写 \(x=L/2\)
的电容虽然能回答前三问，却无法严谨求出液体继续上升时所受的力。

\subsection{五、解析}\label{ux4e94ux89e3ux6790-3}

\begin{figure}[H]
\centering
\includegraphics[width=0.9\linewidth,height=\textheight,keepaspectratio,alt={图4-2　并联分区、表面电荷与液柱受力}]{../figures/solution-04.pdf}
\caption*{图4-2　并联分区、表面电荷与液柱受力}
\end{figure}

图 4-2 把空气区与介质区画成同端电压的并联分区，并标出液柱上升 \(dx\)
时新增的介质体积。下面的
\(C(x)\)、表面自由电荷密度及电场拉力都按图中的面积分区和液柱截面 \(Ld\)
建立。

\needspace{3\baselineskip}
\textbf{（一）充电后断开电池时保存的总电荷}

液体尚未进入时，电容器内为空气，初始电容为

\[
C_0=\frac{\varepsilon_0L^2}{d}. \tag{1}
\]

接在电压为 \(V\) 的电池上充电，故极板自由电荷量的大小为

\[
Q=C_0V
=\frac{\varepsilon_0L^2V}{d}. \tag{2}
\]

断开电池以后，极板与外界没有导电通路，所以在忽略漏电时

\[
\boxed{Q=\text{常量}}. \tag{3}
\]

\needspace{3\baselineskip}
\textbf{（二）第（1）问：液体占据 \(L/2\) 时的电容}

空气部分的极板正对面积为

\[
A_{\rm air}=L\cdot\frac L2=\frac{L^2}{2}.
\]

对应电容为

\[
C_{\rm air}
=\frac{\varepsilon_0A_{\rm air}}d
=\frac{\varepsilon_0L^2}{2d}. \tag{4}
\]

液体部分面积同样为 \(L^2/2\)，介电常数为
\(\varepsilon_r\varepsilon_0\)，故

\[
C_{\rm liq}
=\frac{\varepsilon_r\varepsilon_0L^2}{2d}. \tag{5}
\]

两部分端电压相同，电荷量相加，因此并联：

\[
\begin{aligned}
C
&=C_{\rm air}+C_{\rm liq}\\
&=\frac{\varepsilon_0L^2}{2d}
+\frac{\varepsilon_r\varepsilon_0L^2}{2d}\\
&=\frac{\varepsilon_0L^2}{2d}(1+\varepsilon_r). 
\end{aligned}
\tag{6}
\]

所以

\[
\boxed{C=\frac{\varepsilon_0L^2}{2d}(1+\varepsilon_r)}.
\]

\needspace{3\baselineskip}
\textbf{（三）第（2）问：电压与电场}

断电后电荷量仍为式（2）的 \(Q\)，所以当前电压不是原来的 \(V\)，而是

\[
V'=\frac QC. \tag{7}
\]

代入式（2）、（6）：

\[
\begin{aligned}
V'
&=\frac{\varepsilon_0L^2V/d}
{\varepsilon_0L^2(1+\varepsilon_r)/(2d)}\\
&=\frac{2V}{1+\varepsilon_r}. 
\end{aligned}
\tag{8}
\]

忽略边缘效应，空气区和液体区的两极板间距都为
\(d\)，且端电压相同，所以两区的电场大小相同：

\[
E=\frac{V'}d.
\]

于是

\[
\boxed{
E=\frac{2V}{(1+\varepsilon_r)d}
}. \tag{9}
\]

电场方向由正极板指向负极板。

\needspace{3\baselineskip}
\textbf{（四）第（3）问：金属极板上的自由电荷密度}

空气部分紧邻金属表面处

\[
D_{\rm air}=\varepsilon_0E.
\]

所以正极板空气部分的自由电荷面密度为

\[
\sigma_{\rm air}^{(+)}
=\varepsilon_0E
=\frac{2\varepsilon_0V}{(1+\varepsilon_r)d}. \tag{10}
\]

负极板相应区域符号相反：

\[
\sigma_{\rm air}^{(-)}
=-\frac{2\varepsilon_0V}{(1+\varepsilon_r)d}. \tag{11}
\]

液体部分有

\[
D_{\rm liq}=\varepsilon_r\varepsilon_0E.
\]

因此正极板浸液部分为

\[
\sigma_{\rm liq}^{(+)}
=\varepsilon_r\varepsilon_0E
=\frac{2\varepsilon_r\varepsilon_0V}{(1+\varepsilon_r)d}, \tag{12}
\]

负极板相应为

\[
\sigma_{\rm liq}^{(-)}
=-\frac{2\varepsilon_r\varepsilon_0V}{(1+\varepsilon_r)d}. \tag{13}
\]

所以若只写电荷密度的大小，可概括为

\[
\boxed{
|\sigma_{\rm air}|=\frac{2\varepsilon_0V}{(1+\varepsilon_r)d},
\qquad
|\sigma_{\rm liq}|=\frac{2\varepsilon_r\varepsilon_0V}{(1+\varepsilon_r)d}
}. \tag{14}
\]

检查总电荷：每一区域面积均为 \(L^2/2\)，故正极板总电荷

\[
\begin{aligned}
Q_+
&=\sigma_{\rm air}^{(+)}\frac{L^2}{2}
+\sigma_{\rm liq}^{(+)}\frac{L^2}{2}\\
&=\frac{\varepsilon_0L^2V}{d}=Q,
\end{aligned}
\]

与式（2）一致。

\needspace{3\baselineskip}
\textbf{（五）写出一般浸入长度 \(x\) 时的电容}

令液体在板间上升的长度为 \(x\)。液体部分面积为 \(Lx\)，空气部分面积为
\(L(L-x)\)，所以

\[
\begin{aligned}
C(x)
&=\frac{\varepsilon_0L(L-x)}d
+\frac{\varepsilon_r\varepsilon_0Lx}d\\
&=\frac{\varepsilon_0L}{d}
\left[L+(\varepsilon_r-1)x\right]. 
\end{aligned}
\tag{15}
\]

由

\[
\frac{dC}{dx}
=\frac{\varepsilon_0L}{d}(\varepsilon_r-1)>0
\]

可见液体上升使电容增大。

\needspace{3\baselineskip}
\textbf{（六）固定电荷条件下求电场拉力}

由于 \(Q\) 固定，电场能为

\[
U_E(x)=\frac{Q^2}{2C(x)}. \tag{16}
\]

代入式（15）：

\[
U_E(x)
=\frac{Q^2d}
{2\varepsilon_0L\left[L+(\varepsilon_r-1)x\right]}. \tag{17}
\]

取液体向上进入极板的方向为 \(x\) 正方向，电场对液体的广义力为

\[
F_E=-\frac{dU_E}{dx}. \tag{18}
\]

逐步求导：

\[
\begin{aligned}
\frac{dU_E}{dx}
&=\frac{Q^2d}{2\varepsilon_0L}
\frac{d}{dx}
\left[L+(\varepsilon_r-1)x\right]^{-1}\\
&=-\frac{Q^2d}{2\varepsilon_0L}
(\varepsilon_r-1)
\left[L+(\varepsilon_r-1)x\right]^{-2}.
\end{aligned}
\]

所以

\[
F_E(x)
=\frac{Q^2d(\varepsilon_r-1)}
{2\varepsilon_0L\left[L+(\varepsilon_r-1)x\right]^2}. \tag{19}
\]

稳定状态 \(x=L/2\)，此时

\[
L+(\varepsilon_r-1)\frac L2
=\frac L2(1+\varepsilon_r). \tag{20}
\]

代入式（19）：

\[
\begin{aligned}
F_E
&=\frac{Q^2d(\varepsilon_r-1)}
{2\varepsilon_0L\left[L(1+\varepsilon_r)/2\right]^2}\\
&=\frac{Q^2d(\varepsilon_r-1)}
{2\varepsilon_0L\cdot L^2(1+\varepsilon_r)^2/4}\\
&=\frac{2Q^2d}{\varepsilon_0L^3}
\frac{\varepsilon_r-1}{(1+\varepsilon_r)^2}. 
\end{aligned}
\tag{21}
\]

\needspace{3\baselineskip}
\textbf{（七）第（4）问：与静水压力平衡}

极板内液面比外部液面高 \(h\)，因此板间液体底部相对外部液体存在静水压差

\[
\Delta p=\rho gh. \tag{22}
\]

液柱在水平方向的横截面积为极板宽度 \(L\) 乘板间距 \(d\)：

\[
A_{\rm col}=Ld. \tag{23}
\]

静水压差产生的向下等效力为

\[
F_h=\Delta p\,A_{\rm col}
=\rho ghLd. \tag{24}
\]

稳定时

\[
F_E=F_h. \tag{25}
\]

代入式（21）、（24）：

\[
\frac{2Q^2d}{\varepsilon_0L^3}
\frac{\varepsilon_r-1}{(1+\varepsilon_r)^2}
=\rho ghLd.
\]

约去 \(d\) 并解出 \(h\)：

\[
h
=\frac{2Q^2}{\varepsilon_0\rho gL^4}
\frac{\varepsilon_r-1}{(1+\varepsilon_r)^2}. \tag{26}
\]

由式（2）

\[
Q^2=\frac{\varepsilon_0^2L^4V^2}{d^2}.
\]

代入式（26）：

\[
\begin{aligned}
h
&=\frac{2}{\varepsilon_0\rho gL^4}
\cdot\frac{\varepsilon_0^2L^4V^2}{d^2}
\cdot\frac{\varepsilon_r-1}{(1+\varepsilon_r)^2}\\
&=\frac{2\varepsilon_0V^2}{\rho gd^2}
\frac{\varepsilon_r-1}{(1+\varepsilon_r)^2}. 
\end{aligned}
\tag{27}
\]

所以

\[
\boxed{
h=\frac{2\varepsilon_0V^2}{\rho gd^2}
\frac{\varepsilon_r-1}{(1+\varepsilon_r)^2}
}. \tag{28}
\]

\needspace{3\baselineskip}
\textbf{压力表达式与平衡稳定性}

由 \(Q=CV'\)，式（19）也可写成

\[F_E=\frac12(V')^2\frac{dC}{dx},\qquad
\frac{F_E}{Ld}=\frac12\varepsilon_0(\varepsilon_r-1)E^2.\]

这与液面处电压力跳跃给出的 \(\rho gh\)
相同，可独立核对最终高度差。若电池保持连接，需包括电源做功，结果仍是
\(F_E=V^2C'/2\)，但此时 \(V\) 固定，不能用断电后的
\(V'=2V/(1+\varepsilon_r)\)。

对固定外部液面，液体再上升会使 \(h\) 增大，而固定电荷下 \(F_E(x)\) 随
\(x\)
减小，所以小扰动有恢复趋势。忽略毛细力、边缘场、漏电和介质密度变化；大液池的外液面可视为固定。若计算得到的
\(h\) 不允许液体仅占板间
\(L/2\)，就还需调整极板浸入深度，不能把两个长度直接等同。

\newpage

\section{第 5 题}\label{ux7b2c-5-ux9898}

半径为 \(R\)、质量为 \(M\) 的均匀圆盘起初静止在光滑水平桌面上。质量为
\(m\)、初速度为 \(\bm v_0\) 的质点射向圆盘，击中后嵌入盘缘上的 \(A\)
点。图中 \(C\) 为圆心，直线 \(C'C\) 平行于 \(\bm v_0\)，入射直线与
\(C'C\) 的距离为 \(b\)。点 \(B\) 位于盘缘，且 \(BC\perp C'C\)。已知
\(m\ll M\)。

（1）求使碰后 \(B\) 点相对桌面瞬时静止的 \(b\)。

（2）承接第（1）问，取质点嵌入的时刻为 \(t=0\)，以该时刻 \(B\)
点的位置为固定在桌面上的坐标原点，水平向右为 \(+x\)、竖直向上为
\(+y\)，求 \(A\) 点在时刻 \(t\) 的坐标 \(x(t)\)、\(y(t)\)。

（3）若在盘缘极靠近 \(B\)
点处开孔并以光滑小钉固定，使圆盘可绕小钉转动。嵌入过程中小钉一般承受冲量。求小钉冲量为零时的
\(b\)，并求此时圆盘的角速度。

已知均匀圆盘绕中心轴的转动惯量为 \(I_C=MR^2/2\)。

\begin{figure}[H]
\centering
\includegraphics[width=0.408\linewidth,height=\textheight,keepaspectratio,alt={图5-1　质点入射圆盘的偏心几何}]{../figures/question-05.pdf}
\caption*{图5-1　质点入射圆盘的偏心几何}
\end{figure}

\subsection{一、题型定位}\label{ux4e00ux9898ux578bux5b9aux4f4d-4}

按 \(m\ll M\) 的主阶近似，盘心平动和绕盘心转动可以分开，偏心距约 \(R/2\)
时B点瞬时静止。若保留有限质量修正，就必须改用真实质心，不能只在分母中保留几个
\(m\)。

\subsection{二、知识点}\label{ux4e8cux77e5ux8bc6ux70b9-4}

\begin{enumerate}
\def\labelenumi{\arabic{enumi}.}
\tightlist
\item
  光滑桌面对系统没有水平冲量，水平方向总动量守恒。
\item
  嵌入属于完全非弹性碰撞，机械能不守恒，但动量和关于适当点的角动量仍可守恒。
\item
  自由刚体任一点速度为 \(\bm v_P=\bm V_C+\bm\omega\times\bm r_{P/C}\)。
\item
  质点嵌入后关于圆心的转动惯量近似为 \(I_C'=MR^2/2+mR^2\)。
\item
  点的轨迹可由``平移位置矢量 + 转动位置矢量''直接写出。
\item
  固定轴冲量对轴点的冲量矩为零，因此碰撞前后关于轴点角动量守恒。
\item
  平行轴定理 \(I_B=I_C+MR^2\)。
\end{enumerate}

\subsection{三、物理图像}\label{ux4e09ux7269ux7406ux56feux50cf-4}

图中入射方向取为 \(+y\)。质点从圆心右侧距离 \(b\) 的直线上击中
\(A\)，因而既给系统向上的总动量，也给圆盘一个关于 \(C\)
的角动量。碰后圆心向上匀速运动，圆盘同时转动。位于圆盘左端的 \(B\)
点，其转动速度可以向下；若它恰好抵消圆心向上的平动速度，\(B\)
点瞬时静止。

在第（3）问中，小钉把 \(B\)
固定为真正的转轴。若支点冲量恰为零，碰前质点的线动量就全部成为碰后组合体的总线动量；同时，关于
\(B\) 的角动量又必须守恒。在 \(m/M\to0\) 的主阶近似中，这两个要求给出
\(b\simeq R/2\)；保留有限质量时还要检查水平支点冲量。

\subsection{四、解题入口}\label{ux56dbux89e3ux9898ux5165ux53e3-4}

第（1）问先求圆盘中心的平动速度 \(V_C\) 和角速度 \(\omega\)，再写最左端
\(B\) 的速度 \(v_B=V_C-\omega R\)。题目给出 \(m\ll M\)，因此只保留
\(m/M\) 的主阶，取 \(V_C\simeq mv_0/M\)、\(\omega\simeq2mv_0b/(MR^2)\)。

第（2）问不要从速度分量重新积分。直接写

\[
\bm r_A(t)=\bm r_C(t)+\bm r_{A/C}(t)
\]

即可。第（3）问则换轴到 \(B\)：先求 \(I_B\)，再由
\(I_B\omega=mv_0(R+b)\) 求角速度，并结合冲量为零时的线动量条件
\(\omega R=mv_0/(M+m)\)。

\subsection{五、解析}\label{ux4e94ux89e3ux6790-4}

\begin{figure}[H]
\centering
\includegraphics[width=0.593\linewidth,height=\textheight,keepaspectratio,alt={图5-2　自由圆盘的速度叠加与点 A 的轨迹}]{../figures/solution-05.pdf}
\caption*{图5-2　自由圆盘的速度叠加与点 \(A\) 的轨迹}
\end{figure}

图 5-2 将盘上任一点的速度拆成圆心平动 \(\bm V_C\) 与绕圆心转动
\(\bm\omega\times\bm r\)。第（1）问在 \(B\)
点令这两项抵消，第（2）问则把同一分解积分成 \(A\) 点的参数轨迹。

\needspace{3\baselineskip}
\textbf{原参考的盘心近似计算（有限 \(m\) 项不是受控修正）}

以下保留原参考计算步骤以供对照。其式（9）、（36）、（37）只在
\(m/M\to0\) 的主阶有意义，不能作为精确有限质量答案。

\needspace{3\baselineskip}
\textbf{（一）自由圆盘碰后的平动速度}

取质点入射方向为
\(+y\)。在碰撞的极短时间内，桌面对系统只有竖直方向的支持冲量，而桌面内水平方向没有外冲量。按
\(m\ll M\) 的近似，把圆心 \(C\) 的平动速度记为 \(V_C\)，系统沿 \(y\)
方向动量守恒：

\[
mv_0=(M+m)V_C. \tag{1}
\]

因此

\[
\boxed{V_C=\frac{m}{M+m}v_0}. \tag{2}
\]

碰后圆心沿 \(+y\) 方向匀速运动。

\needspace{3\baselineskip}
\textbf{（二）自由圆盘碰后的角速度}

碰撞前，质点运动直线到圆心 \(C\) 的垂直距离为 \(b\)，所以质点关于 \(C\)
的角动量大小为

\[
L_{C,i}=mv_0b. \tag{3}
\]

碰后质点嵌在盘缘，组合体关于 \(C\) 的转动惯量为

\[
\begin{aligned}
I_C'
&=I_{\rm disk,C}+I_{m,C}\\
&=\frac12MR^2+mR^2\\
&=\left(\frac M2+m\right)R^2. 
\end{aligned}
\tag{4}
\]

关于 \(C\) 的角动量守恒给出

\[
I_C'\omega=mv_0b. \tag{5}
\]

所以

\[
\boxed{
\omega=\frac{mv_0b}{(M/2+m)R^2}
}. \tag{6}
\]

按图示位置，圆盘转动使左端 \(B\) 相对圆心的速度沿 \(-y\) 方向，大小为
\(\omega R\)。

\needspace{3\baselineskip}
\textbf{（三）第（1）问：令 \(B\) 点瞬时静止}

\(B\) 点相对桌面的速度是圆心平动速度与相对圆心转动速度的矢量和。取
\(+y\) 为正：

\[
v_B=V_C-\omega R. \tag{7}
\]

要求 \(B\) 点瞬时静止：

\[
V_C-\omega R=0. \tag{8}
\]

将式（2）、（6）代入：

\[
\frac{m}{M+m}v_0
-\frac{mv_0b}{(M/2+m)R^2}R=0.
\]

约去正数 \(mv_0\)：

\[
\frac1{M+m}
=\frac{b}{(M/2+m)R}.
\]

解出

\[
\boxed{
b=\frac{m+M/2}{m+M}R
}. \tag{9}
\]

由于 \(m\ll M\)，它接近

\[
b\simeq\frac R2.
\]

\needspace{3\baselineskip}
\textbf{（四）第（2）问：确定初始几何角 \(\alpha\)}

以 \(t=0\) 时的 \(B\) 点为坐标原点。由于 \(BC\) 水平向右且长度为
\(R\)，初始圆心位置为

\[
\bm r_C(0)=R\bm i. \tag{10}
\]

圆心以速度 \(V_C\) 沿 \(+y\) 匀速运动，所以

\[
\boxed{
\bm r_C(t)=R\bm i+V_Ct\bm j
}. \tag{11}
\]

质点入射直线与通过 \(C\) 的平行线距离为 \(b\)，因此嵌入点 \(A\)
相对圆心的水平坐标为 \(b\)。设 \(t=0\) 时半径 \(CA\) 位于 \(+x\)
轴下方，与 \(+x\) 轴夹角的大小为 \(\alpha\)，则

\[
R\cos\alpha=b. \tag{12}
\]

所以

\[
\boxed{
\cos\alpha=\frac bR
=\frac{m+M/2}{m+M}
}. \tag{13}
\]

在 \(t=0\)，\(CA\) 的方位角为
\(-\alpha\)。若角速度按图示方向取正，则时刻 \(t\) 的方位角为

\[
\varphi_A(t)=\omega t-\alpha. \tag{14}
\]

于是 \(A\) 相对 \(C\) 的位置矢量为

\[
\boxed{
\bm r_{A/C}(t)
=R\cos(\omega t-\alpha)\bm i
+R\sin(\omega t-\alpha)\bm j
}. \tag{15}
\]

\needspace{3\baselineskip}
\textbf{（五）叠加平动与转动得到轨迹方程}

由

\[
\bm r_A(t)=\bm r_C(t)+\bm r_{A/C}(t), \tag{16}
\]

将式（11）、（15）相加：

\[
\begin{aligned}
\bm r_A(t)
={}&\left[R+R\cos(\omega t-\alpha)\right]\bm i\\
&+\left[V_Ct+R\sin(\omega t-\alpha)\right]\bm j. 
\end{aligned}
\tag{17}
\]

因此

\[
\boxed{
x(t)=R\left[1+\cos(\omega t-\alpha)\right]
} \tag{18}
\]

以及

\[
\boxed{
y(t)=V_Ct+R\sin(\omega t-\alpha)
}. \tag{19}
\]

其中

\[
V_C=\frac{m}{M+m}v_0,
\qquad
\omega=\frac{V_C}{R}
=\frac{mv_0}{(M+m)R}, \tag{20}
\]

最后一个等号使用了第（1）问的 \(B\) 点瞬时静止条件。

检查 \(t=0\)：

\[
x(0)=R(1+\cos\alpha)=R+b,
\]

\[
y(0)=-R\sin\alpha,
\]

正好给出图中 \(A\) 相对 \(B\) 的初始坐标。

\needspace{3\baselineskip}
\textbf{（六）第（3）问：组合体关于 \(B\) 的转动惯量}

现在 \(B\) 处设有固定小钉，碰后圆盘与质点绕 \(B\) 转动。先求均匀圆盘关于
\(B\) 的转动惯量。由平行轴定理：

\[
\begin{aligned}
I_{B,\rm disk}
&=I_{C,\rm disk}+MR^2\\
&=\frac12MR^2+MR^2\\
&=\frac32MR^2. 
\end{aligned}
\tag{21}
\]

再求嵌入点 \(A\) 到 \(B\) 的距离。以 \(C\) 为中间点，\(B\) 在 \(C\) 左侧
\(R\)，而 \(A\) 的水平坐标相对 \(C\) 为 \(b\)，竖直坐标平方为
\(R^2-b^2\)。所以

\[
\begin{aligned}
BA^2
&=(R+b)^2+(R^2-b^2)\\
&=R^2+2Rb+b^2+R^2-b^2\\
&=2R^2+2Rb\\
&=2R(R+b). 
\end{aligned}
\tag{22}
\]

质点关于 \(B\) 的转动惯量为

\[
I_{B,m}=mBA^2=2mR(R+b). \tag{23}
\]

组合体总转动惯量为

\[
\boxed{
I_B=\frac32MR^2+2mR(R+b)
}. \tag{24}
\]

\needspace{3\baselineskip}
\textbf{（七）关于 \(B\) 的角动量守恒}

碰撞过程中，小钉冲量虽然可能很大，但它作用于 \(B\)，对 \(B\)
的冲量矩为零。因此关于 \(B\) 的角动量守恒。

质点入射直线到 \(B\) 的垂直距离为

\[
R+b.
\]

碰前关于 \(B\) 的角动量大小为

\[
L_{B,i}=mv_0(R+b). \tag{25}
\]

碰后为 \(I_B\omega\)，故

\[
I_B\omega=mv_0(R+b). \tag{26}
\]

于是

\[
\omega=\frac{mv_0(R+b)}{I_B}. \tag{27}
\]

\needspace{3\baselineskip}
\textbf{（八）施加``小钉冲量为零''的条件}

设小钉在碰撞过程中给系统的冲量大小为
\(P\)，方向取使系统总动量减少的一向。按题设近似，碰后圆心 \(C\)
的速度大小为

\[
V_C=\omega R. \tag{28}
\]

沿入射方向写冲量---动量方程：

\[
(M+m)V_C-mv_0=-P. \tag{29}
\]

当 \(P=0\) 时：

\[
(M+m)V_C=mv_0. \tag{30}
\]

结合 \(V_C=\omega R\)：

\[
\omega=\frac{mv_0}{(M+m)R}. \tag{31}
\]

另一方面，式（26）给出

\[
I_B\omega=mv_0(R+b).
\]

代入式（31），约去 \(mv_0\)：

\[
I_B\frac1{(M+m)R}=R+b.
\]

即

\[
I_B=(M+m)R(R+b). \tag{32}
\]

再代入式（24）：

\[
\frac32MR^2+2mR(R+b)
=(M+m)R(R+b). \tag{33}
\]

两边除以 \(R\)：

\[
\frac32MR+2m(R+b)
=(M+m)(R+b). \tag{34}
\]

展开：

\[
\frac32MR+2mR+2mb
=MR+Mb+mR+mb.
\]

把含 \(b\) 的项移到一边：

\[
(M-m)b=\left(\frac M2+m\right)R. \tag{35}
\]

所以

\[
\boxed{
b=\frac{m+M/2}{M-m}R
}. \tag{36}
\]

此时角速度由式（31）直接得到：

\[
\boxed{
\omega=\frac{mv_0}{(M+m)R}
}. \tag{37}
\]

由于 \(m\ll M\)，式（36）也近似为
\(b\simeq R/2\)；但它与第（1）问的分母不同，物理约束也不同。

\Needspace{38\baselineskip}

\needspace{3\baselineskip}
\textbf{受控的主阶答案与有限质量问题}

令 \(\epsilon=m/M\ll1\)，一致保留主阶时，两种问题均给出

\[b\simeq R/2,\qquad \omega\simeq\frac{mv_0}{MR},\qquad V_C\simeq\frac{mv_0}{M}.\]

自由盘的A点轨迹在此精度下是

\[x=R[1+\cos(\omega t-\pi/3)],\qquad
 y=R[\omega t+\sin(\omega t-\pi/3)].\]

实际 \(B\) 只是 \(t=0\)
瞬时静止，随后并非固定轴。圆盘加嵌入质点的真实质心为
\(G=C+\eta\boldsymbol a\)，\(\eta=m/(M+m)\)、\(\boldsymbol a=(b,-s)\)、\(s=\sqrt{R^2-b^2}\)。严格的自由碰撞结果为

\[I_G=\frac12MR^2+\frac{Mm}{M+m}R^2,\qquad
\omega=\frac{2mv_0b}{(M+3m)R^2},\qquad \boldsymbol V_G=(0,\eta v_0).\]

因此

\[\boldsymbol v_B=(-\omega\eta s,\ \eta v_0-\omega(R+\eta b)).\]

对实际入射点 \(0<b<R\)，第一分量不为零，所以有限 \(m\)
时不存在B严格静止的偏心距；主阶模型忽略的是这一 \(O(\epsilon^2)\) 速度。

固定B轴时，严格的支点冲量两分量为

\[P_x=m\omega s,\qquad P_y=\omega[MR+m(R+b)]-mv_0,\]

其中 \(\omega=mv_0(R+b)/[3MR^2/2+2mR(R+b)]\)。对于
\(0<b<R\)，\(P_x\ne0\)，所以``小钉冲量为零''也只能在题设近似的主阶意义下成立。原式
\(b=(M/2+m)R/(M-m)\) 忽略了这个横向分量，不是精确零冲量条件。

若需要严格的A轨迹，可写
\(\boldsymbol r_A(t)=\boldsymbol G_0+\boldsymbol V_Gt+(1-\eta)\mathcal R(\omega t)\boldsymbol a\)，其中
\(\boldsymbol G_0=(R,0)+\eta\boldsymbol a\)，\(\mathcal R\)
为平面旋转矩阵。盘心也绕真实质心运动，不再严格匀速直线。

\newpage

\section{第 6 题}\label{ux7b2c-6-ux9898}

一套蚊虫光学探测装置由薄凸透镜和位于其右侧的 CCD
探测器组成。蚊虫发出的特定频段射线经透镜折射后，只要有光线到达
CCD，即认为蚊虫被发现。取透镜平面中心为原点，光轴左方为
\(z>0\)，离轴距离记为 \(R=\sqrt{x^2+y^2}\)。透镜口径为
\(D=10\ \mathrm{cm}\)，焦距为 \(f=10\ \mathrm{cm}\)；CCD 与透镜相距
\(L=15\ \mathrm{cm}\)，CCD 直径为 \(d=4\ \mathrm{cm}\)。

（1）定量描述 \(z>0\) 区域内蚊虫不会被探测到的安全飞行区域，并在
\(R\)---\(z\) 截面内写出边界直线方程。

（2）蚊虫在不被发现的条件下，最小能接近到离光轴多远？

\begin{figure}[H]
\centering
\includegraphics[width=0.82\linewidth,height=\textheight,keepaspectratio,alt={图6-1　透镜、CCD 与空间点的光路几何}]{../figures/question-06.pdf}
\caption*{图6-1　透镜、CCD 与空间点的光路几何}
\end{figure}

\subsection{一、题型定位}\label{ux4e00ux9898ux578bux5b9aux4f4d-5}

把透镜孔径内每一条可能光线映射到CCD平面，检查接收区域是否相交。安全边界由整个光束的包络决定，最近离轴距离是严格安全区的下确界。

\subsection{二、知识点}\label{ux4e8cux77e5ux8bc6ux70b9-5}

\begin{enumerate}
\def\labelenumi{\arabic{enumi}.}
\tightlist
\item
  近轴光线由高度与斜率描述。
\item
  自由传播距离 \(s\) 后，光线高度变为``原高度 + \(s\) 乘斜率''。
\item
  薄透镜使光线斜率改变 \(-\rho/f\)，其中 \(\rho\)
  是光线在透镜平面的高度。
\item
  有限口径条件：透镜处 \(|\rho|\le D/2=5\ \mathrm{cm}\)，CCD 处
  \(|\eta|\le d/2=2\ \mathrm{cm}\)。
\item
  边界光线同时擦过透镜边缘和 CCD 边缘，四种正负组合给出四条直线。
\item
  ``至少存在一条可达光线''可转化为两个一维区间是否相交。
\item
  旋转对称系统只依赖
  \(R=\sqrt{x^2+y^2}\)，二维截面结论绕光轴旋转即为三维区域。
\end{enumerate}

\subsection{三、物理图像}\label{ux4e09ux7269ux7406ux56feux50cf-5}

固定一个空间点 \(P(z,R)\)，让光线在透镜平面的交点高度 \(\rho\) 从
\(-5\ \mathrm{cm}\) 连续变化到 \(+5\ \mathrm{cm}\)。经过透镜后，这束光在
CCD 平面覆盖一个高度区间。如果该区间与 CCD 的 \([-2,2]\ \mathrm{cm}\)
相交，至少有一条光线被接收，蚊虫被发现；若两个区间完全分离，则安全。

安全区的边界发生在两个区间刚好相切，也就是一条极限光线同时通过透镜的上边缘或下边缘，并落在
CCD 的上边缘或下边缘。两种透镜边缘乘两种 CCD
边缘，正好产生四条边界直线。

\subsection{四、解题入口}\label{ux56dbux89e3ux9898ux5165ux53e3-5}

不要先固定``主光线''或只用物像公式。应选透镜面交点高度 \(\rho\)
作为参数，从物点 \(P(z,R)\) 到透镜、再经过薄透镜、最后传播到 CCD，写出
CCD 高度 \(\eta(\rho)\)。这个表达式对 \(\rho\) 是线性的，因此透镜口径在
CCD 上的投影必然是一个区间，边界和最短距离都可直接从区间几何得到。

\subsection{五、解析}\label{ux4e94ux89e3ux6790-5}

\begin{figure}[H]
\centering
\includegraphics[width=0.569\linewidth,height=\textheight,keepaspectratio,alt={图6-2　边缘光线与安全区域边界的构造}]{../figures/solution-06.pdf}
\caption*{图6-2　边缘光线与安全区域边界的构造}
\end{figure}

图 6-2 以透镜交点高度 \(\rho\) 参数化边缘光线，并把透镜口径两端映射到
CCD 上的两个端点。以下写出的 \(\eta(\rho)\)
正是图中这条折线路径的逐段传播结果，包络候选由 \(\rho=\pm a\)、\(a=D/2\)
给出。

\needspace{3\baselineskip}
\textbf{（一）建立单条近轴光线的传输关系}

在包含光轴和物点的子午面内讨论。物点坐标为 \(P(z,R)\)，透镜位于
\(z=0\)。设某条光线在透镜平面的横向高度为 \(\rho\)。

从物点到透镜的轴向距离为 \(z\)，所以入射光线斜率为

\[
\alpha_{\rm in}=\frac{\rho-R}{z}. \tag{1}
\]

在近轴近似下，薄透镜使斜率改变 \(-\rho/f\)，故出射斜率为

\[
\alpha_{\rm out}
=\alpha_{\rm in}-\frac\rho f
=\frac{\rho-R}{z}-\frac\rho f. \tag{2}
\]

CCD 位于透镜右方距离 \(L\) 处。光线从透镜传播到 CCD 后，高度为

\[
\eta=\rho+L\alpha_{\rm out}. \tag{3}
\]

代入式（2）：

\[
\begin{aligned}
\eta
&=\rho+L\left(\frac{\rho-R}{z}-\frac\rho f\right)\\
&=\rho\left(1+\frac Lz-\frac Lf\right)-\frac{LR}{z}. 
\end{aligned}
\tag{4}
\]

使用 \(L=15\ \mathrm{cm}\)、\(f=10\ \mathrm{cm}\)：

\[
1-\frac Lf=1-\frac{15}{10}=-\frac12.
\]

所以

\[
\boxed{
\eta
=\rho\left(\frac{15}{z}-\frac12\right)
-\frac{15R}{z}
}. \tag{5}
\]

式中所有长度统一以厘米计。

\needspace{3\baselineskip}
\textbf{（二）由四条边缘光线求边界直线}

透镜半径为

\[
a=\frac D2=5\ \mathrm{cm},
\]

CCD 半径为

\[
c=\frac d2=2\ \mathrm{cm}. \tag{6}
\]

边界状态取

\[
\rho=\pm5,
\qquad
\eta=\pm2. \tag{7}
\]

把式（5）乘以 \(z\)：

\[
\eta z
=\rho\left(15-\frac z2\right)-15R. \tag{8}
\]

整理出 \(z\)：

\[
z\left(\eta+\frac\rho2\right)=15(\rho-R),
\]

所以

\[
z=\frac{15(\rho-R)}{\eta+\rho/2}. \tag{9}
\]

逐一代入四种组合。

第一种，\(\rho=+5\)、\(\eta=+2\)：

\[
\begin{aligned}
z
&=\frac{15(5-R)}{2+5/2}\\
&=\frac{15(5-R)}{9/2}\\
&=\frac{10}{3}(5-R)\\
&=-\frac{10}{3}R+\frac{50}{3}. 
\end{aligned}
\tag{10}
\]

第二种，\(\rho=-5\)、\(\eta=-2\)：

\[
\begin{aligned}
z
&=\frac{15(-5-R)}{-2-5/2}\\
&=\frac{15(-5-R)}{-9/2}\\
&=\frac{10}{3}(5+R)\\
&=\frac{10}{3}R+\frac{50}{3}. 
\end{aligned}
\tag{11}
\]

第三种，\(\rho=+5\)、\(\eta=-2\)：

\[
\begin{aligned}
z
&=\frac{15(5-R)}{-2+5/2}\\
&=\frac{15(5-R)}{1/2}\\
&=30(5-R)\\
&=-30R+150. 
\end{aligned}
\tag{12}
\]

第四种，\(\rho=-5\)、\(\eta=+2\)：

\[
\begin{aligned}
z
&=\frac{15(-5-R)}{2-5/2}\\
&=\frac{15(-5-R)}{-1/2}\\
&=30(5+R)\\
&=30R+150. 
\end{aligned}
\tag{13}
\]

所以四条边界线为

\[
\boxed{
\begin{aligned}
l_1:&\quad z=-\frac{10}{3}R+\frac{50}{3},\\
l_2:&\quad z= \frac{10}{3}R+\frac{50}{3},\\
l_3:&\quad z=-30R+150,\\
l_4:&\quad z= 30R+150.
\end{aligned}
} \tag{14}
\]

这里 \(R\)
若作为有符号子午坐标，四条线关于光轴成对对称；在三维中应将截面绕 \(z\)
轴旋转。

\needspace{3\baselineskip}
\textbf{（三）用区间相交写出安全区域}

由式（5），当 \(\rho\) 在 \([-5,5]\) 内变化时，CCD 平面上的 \(\eta\)
构成一个区间。该区间中心为

\[
\eta_c=-\frac{15R}{z}, \tag{15}
\]

半宽为

\[
w=5\left|\frac{15}{z}-\frac12\right|. \tag{16}
\]

CCD 有效区间中心为 \(0\)、半宽为 \(2\)。两个区间有交集的条件为

\[
|\eta_c|\le2+w. \tag{17}
\]

由于 \(z>0\) 且取径向距离 \(R\ge0\)：

\[
\frac{15R}{z}
\le2+5\left|\frac{15}{z}-\frac12\right|. \tag{18}
\]

这时蚊虫会被发现。安全区域是严格反向不等式：

\[
R>
\frac z{15}
\left[2+5\left|\frac{15}{z}-\frac12\right|\right]. \tag{19}
\]

根据绝对值分段。

当 \(0<z\le30\ \mathrm{cm}\) 时，\(15/z-1/2\ge0\)：

\[
\begin{aligned}
R
&>\frac z{15}
\left[2+5\left(\frac{15}{z}-\frac12\right)\right]\\
&=\frac{2z}{15}+5-\frac z6\\
&=5-\frac z{30}. 
\end{aligned}
\tag{20}
\]

当 \(z\ge30\ \mathrm{cm}\) 时，\(15/z-1/2\le0\)：

\[
\begin{aligned}
R
&>\frac z{15}
\left[2+5\left(\frac12-\frac{15}{z}\right)\right]\\
&=\frac{2z}{15}+\frac z6-5\\
&=\frac{3z}{10}-5. 
\end{aligned}
\tag{21}
\]

因此在 \(R\ge0\) 半截面内，安全区域可定量写成

\[
\boxed{
\begin{cases}
R>5-\dfrac z{30}, & 0<z\le30\ \mathrm{cm},\\[6pt]
R>\dfrac{3z}{10}-5, & z\ge30\ \mathrm{cm}.
\end{cases}
} \tag{22}
\]

临界线上只有一条边缘光线恰能到达
CCD；若把``接收到任意信号''也算发现，则安全区使用严格大于号。

\needspace{3\baselineskip}
\textbf{（四）第（2）问：求安全区到光轴的最小距离}

安全区边界的径向距离函数为

\[
R_b(z)=
\begin{cases}
5-z/30, & 0<z\le30,\\
3z/10-5, & z\ge30.
\end{cases} \tag{23}
\]

第一段斜率为 \(-1/30\)，函数随 \(z\) 增大而减小；第二段斜率为
\(3/10\)，函数随 \(z\) 增大而增大。因此最小值位于连接点
\(z=30\ \mathrm{cm}\)，代入即得
\(\boxed{R_{\min}=5-30/30=4\ \mathrm{cm}}\)。这表示蚊虫在不被探测到的极限意义下最近可达
\(4\ \mathrm{cm}\)；严格安全时应取 \(R>4\ \mathrm{cm}\)。

\Needspace{30\baselineskip}

\needspace{3\baselineskip}
\textbf{圆孔径的二维判据与边界端点}

轴对称装置也可直接用横向向量表示：

\[\boldsymbol\eta=A(z)\boldsymbol\rho-\frac Lz\boldsymbol R,\qquad
A(z)=1+L/z-L/f.\]

半径 \(a\) 的透镜圆孔径映成CCD平面上半径 \(|A|a\) 的圆盘，圆心离CCD中心
\(LR/z\)。它与半径 \(c\) 的CCD相交当且仅当
\(LR/z\leq c+|A|a\)，因此一维子午截面判据确能推广到三维。

四条边缘线只是候选线，实际安全区内边界是两条包络分段；不能把四条线都当成独立安全边界。\(z=30\ \mathrm{cm}\)
时光束会聚，\(R=4\ \mathrm{cm}\)
恰好像落在CCD边缘，仍被发现。安全点可无限逼近这一值，但没有达到下界的安全点。

这里采用薄透镜近轴模型。口径半径与焦距之比达
\(0.5\)，若用于真实仪器设计，像差和边缘大倾角需要另行处理，题设理想模型的直线包络不能替代实际镜头光线追迹。

\newpage

\section{第 7 题}\label{ux7b2c-7-ux9898}

一个内壁光滑的环形绝热气缸，其截面半径为 \(a\)，中心圆半径为
\(R\)。气缸内等间距放置 \(N\) 个质量均为 \(m\)
的绝热活塞，活塞与内壁无摩擦且不漏气。已知

\[
R\gg\frac RN\gg a,
\]

因此每两个相邻活塞之间的气室可近似为直圆柱。初态每个气室中有
\(\nu\ \mathrm{mol}\) 氮气，温度均为 \(T_0\)。记第 \(i\)
个活塞相对平衡位置的位移为 \(x_i\)，并有 \(|x_i|\ll R/N\)。

（1）求第1、2号活塞之间气体压强相对初态的变化量。

（2）考虑所有活塞都作角频率为 \(\omega\)
的简谐振动、相位随编号线性变化的基频行波模式。相邻活塞相位差
\(\Delta\varphi\ll1\)，可用
\(\cos\Delta\varphi\simeq1-(\Delta\varphi)^2/2\)。又因波速远小于气体声速，每个气室内压强可视为瞬时均匀。求波速的最小可能值。

\begin{figure}[H]
\centering
\includegraphics[width=0.731\linewidth,height=\textheight,keepaspectratio,alt={图7-1　环形绝热气缸与等间距活塞}]{../figures/question-07.pdf}
\caption*{图7-1　环形绝热气缸与等间距活塞}
\end{figure}

\subsection{一、题型定位}\label{ux4e00ux9898ux578bux5b9aux4f4d-6}

相邻活塞的位移差改变气室体积，绝热压强差再推动活塞。先得到离散链方程，再取长波极限；基频表示最低非零频率，不能据此宣称它在全部离散模式中波速最小。

\subsection{二、知识点}\label{ux4e8cux77e5ux8bc6ux70b9-6}

\begin{enumerate}
\def\labelenumi{\arabic{enumi}.}
\tightlist
\item
  环形气缸的周长为 \(2\pi R\)，每个气室的平衡长度为 \(2\pi R/N\)。
\item
  气室截面积为 \(S=\pi a^2\)，平衡体积为 \(V_0=2\pi^2a^2R/N\)。
\item
  理想气体状态方程 \(p_0V_0=\nu\mathcal RT_0\)，其中 \(\mathcal R\)
  为普适气体常量。
\item
  快于热交换但保持准静态均匀的绝热小变化满足
  \(pV^\gamma=\text{常量}\)；氮气近似取 \(\gamma=7/5\)。
\item
  一阶线性化 \((1+\epsilon)^{-\gamma}\simeq1-\gamma\epsilon\)。
\item
  离散耦合方程 \(m\ddot x_n=K_s(x_{n+1}+x_{n-1}-2x_n)\)。
\item
  离散行波试探解 \(x_n=A\cos(\omega t-n\varphi_0)\)。
\item
  环形周期边界条件 \(N\varphi_0=2\pi z\)，其中 \(z\) 为整数。
\end{enumerate}

\subsection{三、物理图像}\label{ux4e09ux7269ux7406ux56feux50cf-6}

若第2号活塞相对第1号活塞沿正方向移动更多，即
\(x_2-x_1>0\)，两者之间的气室变长、体积增大，压强降低。这个较低压强会使两侧活塞受到指向气室内部的恢复趋势。每个活塞同时感受前后两个气室的压力，净力由相邻位移的二阶差分决定。

在基频行波中，各活塞振幅相同，但相位沿编号依次滞后。由于 \(N\)
很大，相邻相位差很小，离散正弦波在局部看起来像连续波。波绕环一周后相位必须增加整数个
\(2\pi\)，基频对应整数 \(z=1\)。

\subsection{四、解题入口}\label{ux56dbux89e3ux9898ux5165ux53e3-6}

先用第1、2号活塞之间的一个气室建立``位移差 \(\to\) 体积变化 \(\to\)
压强变化''的线性关系，并把比例常量记为
\(K\)。随后不再重复热力学推导，只把第 \(n\)
个活塞左右两侧的压强增量写成相邻位移差，得到离散运动方程。最后代入行波解，用三角恒等式和小相位差近似求出色散关系。

\subsection{五、解析}\label{ux4e94ux89e3ux6790-6}

\begin{figure}[H]
\centering
\includegraphics[width=0.9\linewidth,height=\textheight,keepaspectratio,alt={图7-2　相邻气室体积变化与离散行波}]{../figures/solution-07.pdf}
\caption*{图7-2　相邻气室体积变化与离散行波}
\end{figure}

图 7-2 中第 \(n\) 个活塞位移为 \(\xi_n\)，左右气室的体积改变量分别由
\(\xi_n-\xi_{n-1}\) 与 \(\xi_{n+1}-\xi_n\)
决定。下面的离散运动方程将严格按图中这两个位移差写出。

\needspace{3\baselineskip}
\textbf{（一）平衡状态下单个气室的几何量}

环形气缸中心线总长度为

\[
L_{\rm ring}=2\pi R. \tag{1}
\]

\(N\) 个活塞等间距分布，所以相邻活塞之间的平衡弧长为

\[
\ell_0=\frac{2\pi R}{N}. \tag{2}
\]

气缸横截面积为

\[
S=\pi a^2. \tag{3}
\]

单个气室的平衡体积为

\[
\begin{aligned}
V_0
&=S\ell_0\\
&=\pi a^2\cdot\frac{2\pi R}{N}\\
&=\frac{2\pi^2a^2R}{N}. 
\end{aligned}
\tag{4}
\]

由理想气体状态方程：

\[
p_0V_0=\nu\mathcal RT_0. \tag{5}
\]

所以

\[
\begin{aligned}
p_0
&=\frac{\nu\mathcal RT_0}{V_0}\\
&=\frac{N\nu\mathcal RT_0}{2\pi^2a^2R}. 
\end{aligned}
\tag{6}
\]

\needspace{3\baselineskip}
\textbf{（二）第（1）问：第1、2号活塞之间的体积变化}

规定沿编号增大的环向为位移正方向。若第1、2号活塞位移分别为
\(x_1\)、\(x_2\)，两活塞间距的变化量为

\[
\Delta\ell_{12}=x_2-x_1. \tag{7}
\]

因此气室体积变化为

\[
\boxed{
\Delta V_{12}=S(x_2-x_1)
=\pi a^2(x_2-x_1)
}. \tag{8}
\]

当 \(x_2>x_1\)
时气室变长，\(\Delta V_{12}>0\)；这个符号检查与几何直觉一致。

\needspace{3\baselineskip}
\textbf{（三）绝热状态方程的一阶展开}

气缸和活塞均绝热，且运动足够缓慢，使每个气室内部时刻接近平衡。对一个气室有

\[
pV^\gamma=p_0V_0^\gamma, \tag{9}
\]

其中氮气取

\[
\gamma=\frac75. \tag{10}
\]

令 \(V=V_0+\Delta V\)，则

\[
p=p_0\left(\frac{V_0}{V_0+\Delta V}\right)^\gamma.
\]

提出 \(V_0\)：

\[
p=p_0\left(1+\frac{\Delta V}{V_0}\right)^{-\gamma}. \tag{11}
\]

由于 \(|\Delta V|\ll V_0\)，用二项式展开到一阶：

\[
\left(1+\frac{\Delta V}{V_0}\right)^{-\gamma}
\simeq1-\gamma\frac{\Delta V}{V_0}. \tag{12}
\]

故

\[
p\simeq p_0-\gamma p_0\frac{\Delta V}{V_0}. \tag{13}
\]

压强变化量为

\[
\Delta p=p-p_0
=-\gamma p_0\frac{\Delta V}{V_0}. \tag{14}
\]

对第1、2号活塞之间的气室，代入式（8）：

\[
\Delta p_{12}
=-\gamma p_0\frac{\pi a^2}{V_0}(x_2-x_1). \tag{15}
\]

再代入式（4）、（6）。先化简系数：

\[
\begin{aligned}
\frac{p_0\pi a^2}{V_0}
&=\frac{N\nu\mathcal RT_0}{2\pi^2a^2R}
\cdot\frac{\pi a^2}{2\pi^2a^2R/N}\\
&=\frac{N\nu\mathcal RT_0}{2\pi^2a^2R}
\cdot\frac{N}{2\pi R}\\
&=\frac{N^2\nu\mathcal RT_0}{4\pi^3a^2R^2}. 
\end{aligned}
\tag{16}
\]

于是

\[
\Delta p_{12}
=-\frac{\gamma N^2\nu\mathcal RT_0}
{4\pi^3a^2R^2}(x_2-x_1). \tag{17}
\]

取 \(\gamma=7/5\)：

\[
\boxed{
\Delta p_{12}
=-\frac{7N^2\nu\mathcal RT_0}
{20\pi^3a^2R^2}(x_2-x_1)
}. \tag{18}
\]

定义正的压强---位移差系数

\[
K=\frac{\gamma N^2\nu\mathcal RT_0}
{4\pi^3a^2R^2}, \tag{19}
\]

则可简写为

\[
\Delta p_{12}=-K(x_2-x_1). \tag{20}
\]

\needspace{3\baselineskip}
\textbf{（四）建立第 \(n\) 个活塞的离散运动方程}

第 \(n\) 个活塞左侧气室位于活塞 \(n-1\) 与 \(n\) 之间，其压强增量为

\[
\Delta p_{n-1,n}
=-K(x_n-x_{n-1}). \tag{21}
\]

右侧气室位于活塞 \(n\) 与 \(n+1\) 之间，其压强增量为

\[
\Delta p_{n,n+1}
=-K(x_{n+1}-x_n). \tag{22}
\]

平衡压强 \(p_0\) 在左右两侧相互抵消。取正方向为从 \(n\) 指向
\(n+1\)，左侧压强推活塞向正方向，右侧压强推活塞向负方向。因此净力为

\[
F_n=S\left(\Delta p_{n-1,n}-\Delta p_{n,n+1}\right). \tag{23}
\]

代入式（21）、（22）：

\[
\begin{aligned}
F_n
&=S\left[-K(x_n-x_{n-1})+K(x_{n+1}-x_n)\right]\\
&=SK(x_{n+1}+x_{n-1}-2x_n). 
\end{aligned}
\tag{24}
\]

由牛顿第二定律：

\[
\boxed{
m\ddot x_n
=\pi a^2K(x_{n+1}+x_{n-1}-2x_n)
}. \tag{25}
\]

若某活塞相对两邻居向正方向凸出，即
\(x_n>(x_{n-1}+x_{n+1})/2\)，右端括号为负，净力指向负方向，确为恢复力。

\needspace{3\baselineskip}
\textbf{（五）代入离散行波试探解}

设

\[
x_n=A\cos(\omega t-n\varphi_0), \tag{26}
\]

其中 \(\varphi_0\) 为相邻活塞的相位差。则

\[
\ddot x_n=-\omega^2x_n. \tag{27}
\]

相邻两项为

\[
x_{n+1}=A\cos(\omega t-n\varphi_0-\varphi_0),
\]

\[
x_{n-1}=A\cos(\omega t-n\varphi_0+\varphi_0).
\]

利用恒等式

\[
\cos(\alpha-\beta)+\cos(\alpha+\beta)
=2\cos\alpha\cos\beta,
\]

得到

\[
x_{n+1}+x_{n-1}=2x_n\cos\varphi_0. \tag{28}
\]

因此

\[
x_{n+1}+x_{n-1}-2x_n
=2x_n(\cos\varphi_0-1). \tag{29}
\]

把式（27）、（29）代入式（25）：

\[
-m\omega^2x_n
=2\pi a^2K(\cos\varphi_0-1)x_n. \tag{30}
\]

约去非零振幅对应的 \(x_n\)：

\[
m\omega^2
=2\pi a^2K(1-\cos\varphi_0). \tag{31}
\]

这就是精确的离散色散关系。

\needspace{3\baselineskip}
\textbf{（六）长波近似下的角频率}

相邻相位差很小，使用

\[
\cos\varphi_0\simeq1-\frac{\varphi_0^2}{2}. \tag{32}
\]

所以

\[
1-\cos\varphi_0\simeq\frac{\varphi_0^2}{2}. \tag{33}
\]

代入式（31）：

\[
m\omega^2
\simeq2\pi a^2K\cdot\frac{\varphi_0^2}{2}
=\pi a^2K\varphi_0^2. \tag{34}
\]

因此

\[
\boxed{
\omega=\sqrt{\frac{\pi a^2K}{m}}\,\varphi_0
}. \tag{35}
\]

\needspace{3\baselineskip}
\textbf{（七）周期边界条件与波速}

波绕环一周后必须回到同一相位，因此

\[
N\varphi_0=2\pi z,
\qquad z=0,1,2,\ldots \tag{36}
\]

\(z=0\)
对应所有活塞同相整体平移，不引起气室体积变化，频率为零。最低的非零振动模式是

\[
z=1,
\qquad
\varphi_0=\frac{2\pi}{N}. \tag{37}
\]

相邻活塞的空间距离为
\(\ell_0=2\pi R/N\)。波从一个活塞传播到下一个活塞时，相位延迟为
\(\varphi_0\)，对应时间延迟

\[
\Delta t=\frac{\varphi_0}{\omega}. \tag{38}
\]

所以相速度为

\[
\begin{aligned}
v_{\rm w}
&=\frac{\ell_0}{\Delta t}\\
&=\frac{2\pi R/N}{\varphi_0/\omega}\\
&=\frac{2\pi R}{N}\frac{\omega}{\varphi_0}. 
\end{aligned}
\tag{39}
\]

代入式（35）：

\[
v_{\rm w}
=\frac{2\pi R}{N}
\sqrt{\frac{\pi a^2K}{m}}. \tag{40}
\]

再代入式（19）：

\[
\begin{aligned}
v_{\rm w}
&=\frac{2\pi R}{N}
\sqrt{
\frac{\pi a^2}{m}
\cdot
\frac{\gamma N^2\nu\mathcal RT_0}
{4\pi^3a^2R^2}
}\\
&=\frac{2\pi R}{N}
\sqrt{
\frac{\gamma N^2\nu\mathcal RT_0}
{4\pi^2mR^2}
}\\
&=\frac{2\pi R}{N}
\cdot\frac{N}{2\pi R}
\sqrt{\frac{\gamma\nu\mathcal RT_0}{m}}\\
&=\sqrt{\frac{\gamma\nu\mathcal RT_0}{m}}. 
\end{aligned}
\tag{41}
\]

取氮气 \(\gamma=7/5\)：

\[
\boxed{
v_{\min}=\sqrt{\frac{7\nu\mathcal RT_0}{5m}}
}. \tag{42}
\]

在题给长波近似下，各低阶模式共享此波速。基频 \(z=1\)
对应最低非零频率，且其波长为周长，此时有

\[
v_{\min}=R\omega_1. \tag{43}
\]

\needspace{3\baselineskip}
\textbf{有限N色散与``最小''的含义}

令
\(\ell_0=2\pi R/N\)、\(c_0=\sqrt{\gamma\nu\mathcal RT_0/m}\)，精确离散关系为

\[\omega(k)=\frac{2c_0}{\ell_0}|\sin(k\ell_0/2)|.\]

基频 \(k_1=1/R\) 的相速度为

\[c_1=c_0\frac{\sin(\pi/N)}{\pi/N}
=c_0[1-\pi^2/(6N^2)+O(N^{-4})].\]

它趋于正文的 \(c_0\)。若纳入整个布里渊区，\(\sin x/x\) 随 \(x\)
增大而减小，更高阶模式反而有更小相速度。因此原题``波速最小''只有在规定的基频、长波估算语境下解释，不能扩展为所有模式的全局极小值。

波速远小于气体声速要求每室的活塞惯性远大于气体惯性，粗略为
\(m\gg\nu M_{\rm mol}\)；它保证室内迅速均压，却不把绝热过程改成等温过程。

\newpage

\section{第 8 题}\label{ux7b2c-8-ux9898}

一束缪子在均匀磁场中沿圆形轨道运动，忽略电磁辐射造成的能量损失。缪子质量为

\[
m_\mu=1.88\times10^{-28}\ \mathrm{kg},
\]

电荷量为

\[
q_\mu=-1.602\times10^{-19}\ \mathrm C,
\]

半衰期为

\[
t_{1/2}=1.523\ \mu\mathrm s.
\]

（1）当缪子速度远小于光速时，若每运动一周恰有一半缪子衰变，求磁感应强度
\(B\)。

（2）保持同一磁场而提高缪子速度，使其不再远小于光速。问缪子每运动一周的衰变比例增大、减小还是保持不变？

\begin{figure}[H]
\centering
\includegraphics[width=0.515\linewidth,height=\textheight,keepaspectratio,alt={图8-1　缪子在均匀磁场中的圆周运动}]{../figures/question-08.pdf}
\caption*{图8-1　缪子在均匀磁场中的圆周运动}
\end{figure}

\subsection{一、题型定位}\label{ux4e00ux9898ux578bux5b9aux4f4d-7}

每绕一周积累的固有时由磁场和粒子静质量决定。提速后实验室的回旋周期与衰变寿命同乘
\(\gamma\)，所以每周存活比例保持不变。

\subsection{二、知识点}\label{ux4e8cux77e5ux8bc6ux70b9-7}

\begin{enumerate}
\def\labelenumi{\arabic{enumi}.}
\tightlist
\item
  非相对论磁洛伦兹力提供向心力：\(|q|vB=mv^2/r\)。
\item
  非相对论回旋角频率 \(\omega_c=|q|B/m\)，周期
  \(T_0=2\pi m/(|q|B)\)，与速度无关。
\item
  半衰期定义：经过一个半衰期，存活粒子数降为原来的一半。
\item
  相对论动量 \(\bm p=\gamma m\bm v\)；磁场不做功，速度大小与 \(\gamma\)
  保持不变。
\item
  相对论回旋角频率 \(\omega_c=|q|B/(\gamma m)\)，实验室周期
  \(T=\gamma T_0\)。
\item
  运动缪子在实验室中的半衰期为 \(t_{1/2}^{\rm lab}=\gamma t_{1/2}\)。
\item
  衰变概率取决于沿粒子世界线积累的固有时。
\end{enumerate}

\subsection{三、物理图像}\label{ux4e09ux7269ux7406ux56feux50cf-7}

均匀磁场始终给缪子一个垂直于速度的力，只改变速度方向，不改变速度大小。因此缪子在每一周中保持同一个洛伦兹因子
\(\gamma\)。

提高速度后，实验室观察到两个过程同时变慢：一方面缪子的动量增大，同一磁场转弯能力变弱，绕一周需要更长时间；另一方面缪子的内部衰变钟也按完全相同的
\(\gamma\) 倍变慢。因此一周所对应的固有时间不变，每周衰变比例也不变。

\subsection{四、解题入口}\label{ux56dbux89e3ux9898ux5165ux53e3-7}

第一问直接令非相对论周期 \(T_0\)
等于半衰期。第二问不要只说``寿命变长，所以衰变减少''，也不要只说``周期变长，所以衰变增加''；必须把实验室一周时间
\(T=\gamma T_0\) 与实验室半衰期 \(\gamma t_{1/2}\)
取比值，才能判断每周衰变比例。

\subsection{五、解析}\label{ux4e94ux89e3ux6790-7}

\begin{figure}[H]
\centering
\includegraphics[width=0.9\linewidth,height=\textheight,keepaspectratio,alt={图8-2　回旋周期与衰变时间的共同相对论伸长}]{../figures/solution-08.pdf}
\caption*{图8-2　回旋周期与衰变时间的共同相对论伸长}
\end{figure}

图 8-2
把实验室中的回旋一周时间与衰变计时放在同一时间轴上；两者都乘同一个洛伦兹因子
\(\gamma\)。后文将直接比较 \(T=\gamma T_0\) 与
\(t_{1/2}^{\rm lab}=\gamma t_{1/2}\)，避免只考虑其中一种时间膨胀。

\needspace{3\baselineskip}
\textbf{（一）非相对论回旋周期}

缪子电荷为负，磁场决定其转向，但求周期只需使用电荷量绝对值
\(|q_\mu|\)。磁洛伦兹力大小为

\[
F_B=|q_\mu|vB. \tag{1}
\]

它提供匀速圆周运动所需向心力：

\[
|q_\mu|vB=\frac{m_\mu v^2}{r}. \tag{2}
\]

约去 \(v\)：

\[
|q_\mu|B=\frac{m_\mu v}{r}. \tag{3}
\]

圆周运动周期为

\[
T_0=\frac{2\pi r}{v}. \tag{4}
\]

由式（3）

\[
\frac rv=\frac{m_\mu}{|q_\mu|B}.
\]

代入式（4）：

\[
\boxed{
T_0=\frac{2\pi m_\mu}{|q_\mu|B}
}. \tag{5}
\]

速度 \(v\) 与轨道半径 \(r\) 同时消去，所以非相对论回旋周期与速度无关。

\needspace{3\baselineskip}
\textbf{（二）第（1）问：周期等于半衰期}

``每运动一周有一半缪子衰变''表示一周后存活数为初始数的一半。因此一周所用时间正好等于半衰期：

\[
T_0=t_{1/2}. \tag{6}
\]

由式（5）：

\[
B=\frac{2\pi m_\mu}{|q_\mu|t_{1/2}}. \tag{7}
\]

把半衰期换成秒：

\[
t_{1/2}=1.523\times10^{-6}\ \mathrm s. \tag{8}
\]

代入数值：

\[
\begin{aligned}
B
&=\frac{2\pi\times1.88\times10^{-28}}
{1.602\times10^{-19}\times1.523\times10^{-6}}
\ \mathrm T\\
&=\frac{1.18124\times10^{-27}}
{2.43985\times10^{-25}}
\ \mathrm T\\
&=4.84\times10^{-3}\ \mathrm T. 
\end{aligned}
\tag{9}
\]

所以

\[
\boxed{B\simeq4.84\times10^{-3}\ \mathrm T}
\]

即约为

\[
\boxed{B\simeq4.84\ \mathrm{mT}}. \tag{10}
\]

\needspace{3\baselineskip}
\textbf{（三）相对论回旋周期}

提高速度后，缪子动量为

\[
\bm p=\gamma m_\mu\bm v,
\qquad
\gamma=\frac1{\sqrt{1-v^2/c^2}}. \tag{11}
\]

由于磁场不做功，\(v\) 的大小不变，所以 \(\gamma\)
在运动中也不变。垂直磁场的动量方向以角速度旋转。用相对论形式的向心关系：

\[
|q_\mu|vB=\frac{\gamma m_\mu v^2}{r}. \tag{12}
\]

约去 \(v\)：

\[
r=\frac{\gamma m_\mu v}{|q_\mu|B}. \tag{13}
\]

实验室中一周周期为

\[
\begin{aligned}
T
&=\frac{2\pi r}{v}\\
&=\frac{2\pi}{v}
\cdot\frac{\gamma m_\mu v}{|q_\mu|B}\\
&=\frac{2\pi\gamma m_\mu}{|q_\mu|B}. 
\end{aligned}
\tag{14}
\]

与式（5）比较：

\[
\boxed{T=\gamma T_0}. \tag{15}
\]

所以速度越接近光速，实验室看到的回旋周期越长。

\needspace{3\baselineskip}
\textbf{（四）运动缪子的半衰期}

题给半衰期 \(t_{1/2}\)
是缪子静止系中的固有半衰期。运动缪子的内部时钟相对实验室变慢，因此实验室测得的半衰期为

\[
\boxed{
t_{1/2}^{\rm lab}=\gamma t_{1/2}
}. \tag{16}
\]

在本题所用同一磁场中，非相对论周期已由第一问调到

\[
T_0=t_{1/2}. \tag{17}
\]

所以相对论情况下

\[
\frac{T}{t_{1/2}^{\rm lab}}
=\frac{\gamma T_0}{\gamma t_{1/2}}
=\frac{T_0}{t_{1/2}}
=1. \tag{18}
\]

经过一个实验室半衰期，存活比例仍为 \(1/2\)。因此每运动一周：

\[
\boxed{
\text{仍有一半缪子衰变，衰变比例保持不变}
}. \tag{19}
\]

\needspace{3\baselineskip}
\textbf{（五）用固有时再次核验}

缪子在实验室运动一周所积累的固有时为

\[
\Delta\tau=\frac T\gamma. \tag{20}
\]

代入 \(T=\gamma T_0\)：

\[
\Delta\tau=T_0. \tag{21}
\]

而第一问有
\(T_0=t_{1/2}\)，所以不论速度多大，只要磁场相同且能量不损失，每一周在缪子自身时钟上都恰好经过一个半衰期。这与式（19）完全一致。

\needspace{3\baselineskip}
\textbf{用周数写出衰变规律}

对任意同一恒定磁场，忽略辐射时一周固有时为
\(\tau_B=2\pi m_\mu/(|q_\mu|B)\)。存活比例可写成

\[\frac{N_n}{N_0}=2^{-n\tau_B/t_{1/2}}.\]

题设调成 \(\tau_B=t_{1/2}\)，故第 \(n\) 周后剩
\(2^{-n}\)。磁场加倍时，一周固有时减半，每周存活比例变成
\(1/\sqrt2\)，衰变比例为
\(1-1/\sqrt2\)。若要求每周只剩四分之一，则一周需两个半衰期，磁场应减半。

半衰期是静止系给定的衰变参数。圆周运动有加速度，但累计固有时仍可沿世界线积分；不应把整个轨道套成一个全局匀速惯性系。
\end{document}
